% ==================================================
% Fixed-Point Causality:
% Physics, Consciousness, and Computation in the Relativistic Scalar–Vector Plenum
% ==================================================

\documentclass[12pt,a4paper]{article}

% --- Encoding & fonts
\usepackage[utf8]{inputenc}
\usepackage[T1]{fontenc}
\usepackage{lmodern}
\usepackage{microtype}

% --- Math & symbols
\usepackage{amsmath,amssymb,amsthm,mathtools,bm}

% --- Layout & lists
\usepackage{geometry}
\geometry{margin=1in}
\usepackage{enumitem}
\usepackage{booktabs}
\usepackage{multicol}

% --- Hyperlinks
\usepackage{hyperref}
\hypersetup{
  colorlinks=true,
  linkcolor=blue,
  citecolor=blue,
  urlcolor=cyan,
  pdfauthor={},
  pdftitle={Fixed-Point Causality: Physics, Consciousness, and Computation in the Relativistic Scalar--Vector Plenum}
}

% --- Citations (author-year)
\usepackage[round,authoryear]{natbib}
% Author–year punctuation similar to many philosophy/physics journals
\bibpunct{(}{)}{;}{a}{}{,}

% --- Custom macros for RSVP framework
\newcommand{\Ph}{\Phi}                % scalar potential
\newcommand{\Vv}{\mathbf{v}}          % vector flow
\newcommand{\Ss}{S}                   % entropy density / functional
\newcommand{\RSVP}{RSVP}              % framework name
\newcommand{\CLIO}{CLIO}              % cognitive module
\newcommand{\TARTAN}{TARTAN}          % semantic tiling framework
\newcommand{\lag}{\mathcal{L}}        % Lagrangian density
\newcommand{\action}{\mathcal{A}}     % action
\newcommand{\divv}{\nabla\!\cdot\!}   % divergence
\newcommand{\grad}{\nabla}            % gradient
\newcommand{\ddt}{\tfrac{d}{dt}}      % time derivative
\newcommand{\Eval}{\mathcal{E}}       % generic evaluator
\newcommand{\State}{\Psi}             % global state
\newcommand{\F}{\mathcal{F}}          % fixed-point evaluator
\newcommand{\LambdaOp}{\Lambda}       % lazy continuation operator

% --- Theorem-like environments (kept minimal for monograph style)
\theoremstyle{plain}
\newtheorem{proposition}{Proposition}
\newtheorem{definition}{Definition}

% ==================================================
\title{\Huge\bfseries Fixed--Point Causality\\[0.35em]
\Large Physics, Consciousness, and Computation in the Relativistic Scalar--Vector Plenum}
\author{Flyxion}
\date{October 2025}

\begin{document}
\maketitle

\begin{abstract}
Modern theories of physics, cognition, and computation face a common impasse: evaluation appears to require an external frame. The \emph{Relativistic Scalar--Vector Plenum} (\RSVP) advances a remedy by treating energy, information, and meaning as co-evolving fields whose \emph{stability}---rather than their motion through a background---constitutes reality. Within this setting, causality is not a chain of prior events but a condition of \emph{equilibrium}: an event occurs when further evaluation no longer alters its outcome. We formalize this as \emph{Fixed--Point Causality} (FPC): processes persist insofar as they are invariant under their own transformation.

FPC dissolves the regress of self-reference by replacing recursive repetition with evaluative closure. The same operator grammar governs (i) a physical layer, where scalar potential (\(\Ph\)), vector flow (\(\Vv\)), and entropy (\(\Ss\)) mutually relax; (ii) a cognitive layer (\CLIO), where predictive inference halts when expected and observed information coincide; and (iii) a computational--semantic layer, where lazy continuation and categorical tiling (\TARTAN) converge on idempotent meanings. The result links thermodynamic, cognitive, and algorithmic halting via one invariant: \(\F[\State]=\State \iff \ddt \Ss=0\). We develop the mathematics, empirical implications, and conceptual unification, showing how “meaning’’ can be identified with the \emph{stable coincidence} of observation and existence.
\end{abstract}

\tableofcontents
\newpage

% ==================================================
% Chapter 1
\section{Introduction}
\label{sec:introduction}

\subsection{Historical context: from expansion and representation to evaluation}
The recurrent difficulty in theories of the universe and the mind is not merely what exists, but \emph{how} existence is stabilized. In standard cosmology, structure is often explained via metric expansion, yet the explanatory burden is displaced onto the origin and persistence of the metric itself. Thermodynamic approaches have shown that gravitational dynamics can be read off equilibrium conditions \citep{jacobson1995thermodynamics,verlinde2011origin}, suggesting that geometry may be a bookkeeping of information flow rather than an a priori container. In cognitive science, predictive-processing and active-inference accounts model perception as minimizing expected surprise or free energy \citep{friston2023active,clark2023brains,seth2022predictive}, but must still specify what makes an inference \emph{complete}. In computation, recursion provides expressive power at the price of regress: a call that calls itself presupposes an external evaluator or base case to terminate evaluation.

The \RSVP\ framework reframes these issues by treating the cosmos as a \emph{field of evaluation}. The fundamental explanandum is stability: what persists are those configurations for which further evaluation produces no net informational change. On this reading, dynamics are not primitive movers but the visible relaxation of gradients toward equilibrium. The thesis aligns with thermodynamic derivations of gravitational field equations \citep{jacobson1995thermodynamics} and entropic readings of force \citep{verlinde2011origin}, yet departs from expansionist narratives by placing the evaluative condition itself---rather than an expanding metric---at the heart of causation.

\subsection{From recursion to continuation: the grammar of closure}
Recursion is a syntactic device: \(f(x)=f(f(x))\) defines self-application but not \emph{when} to stop. Continuation, by contrast, is semantic: it advances only insofar as evaluation has not yet stabilized. We therefore elevate a \emph{fixed-point} criterion,
\begin{equation}
\label{eq:fpc-law}
\F[\State]=\State
\qquad\Longleftrightarrow\qquad
\ddt \Ss=0,
\end{equation}
as the universal halting condition. In this view, an event is a local attainment of informational equilibrium, not a tick in an external time. The move generalizes the stationarity logic of variational principles to any evaluative field: just as stationary action identifies physical trajectories, fixed points of \(\F\) identify completed evaluations. The practical consequence is an ontology that unifies physics, cognition, and computation under a single \emph{grammar of closure}.

\subsection{Outline and contributions}
This monograph develops Fixed--Point Causality as follows. Chapter~\ref{sec:fpc} formalizes the evaluator \(\F\), the lazy continuation operator \(\LambdaOp\), and the contrast with recursive, dynamical, and variational paradigms. Chapter~\ref{sec:rsvp-field} derives the \RSVP\ field equations for \(\Ph\), \(\Vv\), and \(\Ss\) from an action with entropic constraints, reading spacetime as emergent equilibrium structure \citep{jacobson1995thermodynamics,verlinde2011origin,carney2023gravity}. Chapter~\ref{sec:origins} clarifies the methodological avoidance of explicit recursion and introduces \emph{environmental recursion} as re-entry of evaluation through changing context (Bateson’s cybernetics and Rosen’s anticipatory systems provide useful antecedents; \citealt{bateson1972steps,rosen1985anticipatory}). Chapter~\ref{sec:clio} defines \CLIO\ as cognitive fixed-point search, relating predictive closure to measurable entropy-flow signatures \citep{friston2023active,seth2022predictive}. Chapter~\ref{sec:computational} reinterprets fixed-point combinators and halting as thermodynamic sufficiency \citep{turing1936computable,kleene1952recursion,barendregt1984lambda,goodhart1975problems}. Chapter~\ref{sec:comparanda} situates FPC relative to Platonic invariance, critiques behaviorist RLHF \citep{christiano2017deep,leike2018scalable,gabriel2020alignment}, and frames computation as geometric vector transformation on complex manifolds with attractor selection. Chapter~\ref{sec:tartan} formulates \TARTAN\ as semantic tiling with homotopy colimits \citep{lawvere2009conceptual,baez2011physics,spivak2014database}. Chapter~\ref{sec:cogphys} articulates the cognitive--physical isomorphism and introduces \emph{viviception}. Chapter~\ref{sec:implications} collects theoretical payoffs, empirical predictions, and applications. Chapter~\ref{sec:conclusion} concludes. Appendices provide derived geometry of coarse-graining, numerical unistochasticity, a historical note on rotational ontology, and the proximal fixed-point operator.

% ==================================================
% Chapter 2
\section{Fixed--Point Causality (FPC)}
\label{sec:fpc}

\subsection{Definition: causality as invariance under evaluation}
We define causality as persistence under self-transformation. Let \(\State=(\Ph,\Vv,\Ss)\) denote the global evaluative state of the plenum. An \emph{evaluator} \(\F\) acts on states to test whether further change is informationally warranted. Causal completion occurs when
\begin{equation}
\F[\State]=\State,
\qquad\text{equivalently}\qquad
\ddt \Ss=0,
\end{equation}
so that the entropy gradient driving evaluation has vanished. Unlike recursion, which requires a meta-level call to iterate itself, \(\F\) is reflexive without regress: it does not execute its argument repeatedly but checks for invariance of informational content. In this sense, being is \emph{invariance under evaluation}.

\subsection{Lazy continuation and proximal stationarity}
To capture the dynamics of approach to equilibrium, introduce a \emph{lazy continuation} operator \(\LambdaOp\) that updates only unresolved gradients:
\begin{equation}
\label{eq:lazy}
\State_{t+1}=\LambdaOp(\State_t),
\qquad
\State_*=\LambdaOp(\State_*).
\end{equation}
A convenient realization is proximal iteration on a local Lagrangian density \(\lag\),
\begin{equation}
\label{eq:prox}
\State_{t+1}
= \operatorname{prox}_{\tau\lag}(\State_t)
:= \arg\min_{\State'}\Big[\,\lag(\State')+\tfrac{1}{2\tau}\|\State'-\State_t\|^2\,\Big],
\end{equation}
with relaxation parameter \(\tau>0\). At a fixed point \(\State_{t+1}=\State_t\), the Euler condition implies \(\nabla \lag(\State_*)=0\). Equation~\eqref{eq:prox} exhibits the structural equivalence between variational stationarity and evaluative closure: both halt when no further descent is possible. Thermodynamically, this is \(\ddt \Ss=0\); computationally, it is halting by sufficiency rather than by fiat.

\subsection{Ontological and epistemic reading}
FPC reorients ontology from substance to stability and epistemology from representation to evaluation. To exist, on this account, is to remain self-identical under the transformations that define existence. To know is not to mirror being, but to operate so that being becomes self-consistent. This resonates with Dewey’s logic of inquiry as stabilization of an indeterminate situation \citep{dewey1938logic} and with the structural ambitions of categorical physics \citep{baez2011physics}. In FPC, truth is not correspondence to an external template but the attainment of idempotence under \(\F\).

\subsection{Comparative paradigms}
A comparative synopsis clarifies the distinctiveness of FPC:

\begin{center}
\begin{tabular}{@{}lllll@{}}
\toprule
\textbf{Paradigm} & \textbf{Operator Form} & \textbf{Evaluator} & \textbf{Stability} & \textbf{Causal Type} \\
\midrule
Recursive   & \(f(f(x))\)            & External meta-call   & Divergent/undecidable & Iterative \\
Dynamical   & \(\dot{x}=F(x)\)       & Temporal derivative  & Oscillatory/attractor & Temporal \\
Variational & \(\delta \mathcal{S}=0\)& Functional extremizer& Extremal              & Stationary \\
\textbf{Fixed-Point} & \(\F(x)=x\)     & Internal evaluator   & \textbf{Stable}       & \textbf{Evaluative} \\
\bottomrule
\end{tabular}
\end{center}

Recursive accounts require an external controller to halt; dynamical accounts presuppose a background time; variational accounts posit an action to be extremized. Fixed-point causality alone renders evaluation \emph{self-contained}: halting is explained by the vanishing of the very gradients that propel change. This is the grammar in which physics, cognition, and computation can be read as different \emph{projections} of a single operation.

\subsection{Remarks on computation and undecidability}
Within standard computability, the halting problem is undecidable because termination is framed syntactically: no general algorithm decides from a program text whether execution stops \citep{turing1936computable}. FPC reframes termination as a \emph{thermodynamic} property: evaluation halts when the entropy gradient is zero. While this does not solve the classical undecidability result at the level of Turing descriptions, it provides a physically meaningful criterion of sufficiency for real evaluators. As \citet{goodhart1975problems} warned for optimization, driving a proxy beyond its regime distorts the target; FPC treats this distortion as informational noise produced after equilibrium has been achieved. The prescription is simple: evaluate until gradients vanish, then stop.

% ==================================================
% Chapter 3
\section{The \RSVP\ Field Theory under Fixed--Point Causality}
\label{sec:rsvp-field}

\subsection{Motivation and structure of the plenum}
The Relativistic Scalar--Vector Plenum (\RSVP) is conceived as a continuous evaluative medium that replaces the expanding metric of general relativity with the dynamics of informational relaxation.
Instead of a background spacetime upon which matter acts, \RSVP\ posits an intertwined set of scalar, vector, and entropic fields whose fixed-point relations generate both geometry and phenomenology.
The essential postulate is that causal structure \emph{emerges} when evaluation halts locally:
a region of the plenum is said to exist when its internal transformations yield no net informational change.

Let the global state of the plenum be
\(\State = (\Ph,\Vv,\Ss)\),
where \(\Ph\) denotes scalar potential, \(\Vv\) denotes vector flow, and \(\Ss\) is local entropy density.
Each component interacts through coupled conservation and relaxation laws that ensure the total informational flux of the system approaches equilibrium.
The field theory is constructed to make Eq.~\eqref{eq:fpc-law} explicit:
\[
\F[\Ph,\Vv,\Ss] = (\Ph,\Vv,\Ss).
\]
This definition anchors all further derivations: the physical universe corresponds to those configurations that are invariant under the evaluator \(\F\).

\subsection{Governing equations}
At macroscopic scales, the scalar and vector fields obey coupled evolution equations reminiscent of fluid dynamics and gauge theory but interpreted thermodynamically:
\begin{align}
\partial_t \Ph &= -\grad\!\cdot(\Ph \Vv) + \xi_\Ph(t), \label{eq:phi-evo}\\
\partial_t \Vv &= -\grad \Ss - \Ph \grad\!\times\!\Vv + \xi_\Vv(t), \label{eq:v-evo}\\
\partial_t \Ss &= -\divv(\Ss\Vv) + \kappa \nabla^2 \Ss + \xi_\Ss(t), \label{eq:s-evo}
\end{align}
where \(\xi_i(t)\) are stochastic noise sources representing unresolved microstructure and \(\kappa\) is a diffusion coefficient that governs the rate of entropic smoothing.
Equations~\eqref{eq:phi-evo}--\eqref{eq:s-evo} express that scalar potential drives vector flow, vector flow transports entropy, and entropy gradients in turn relax the scalar potential---a closed triadic system.
In equilibrium, the time derivatives vanish and each equation reduces to a fixed-point constraint.
The steady state of these coupled PDEs defines the macroscopic geometry that appears to observers as “spacetime.”

\subsection{Variational formulation}
To capture the dynamics of convergence more elegantly, we introduce a Lagrangian density
\begin{equation}
\lag(\Ph,\Vv,\Ss,\partial_\mu\Ph,\partial_\mu\Vv)
 = \frac{1}{2}(\partial_\mu \Ph)(\partial^\mu \Ph)
   + \frac{1}{2}(\grad\!\times\!\Vv)^2
   - V(\Ph,\Vv,\Ss),
\end{equation}
where \(V\) encodes the local entropic potential of the plenum.
The corresponding action functional is
\begin{equation}
\label{eq:action}
\action = \int \lag\, d^4x.
\end{equation}
Stationarity of \(\action\) yields Euler–Lagrange equations that coincide with the relaxation equations above when the potential satisfies
\(\partial V/\partial \Ph = \divv(\Ph \Vv)\)
and
\(\partial V/\partial \Vv = \grad \Ss\).
Thus, the physical trajectories of the fields correspond to extrema of \(\action\), which themselves are fixed points under the evaluator \(\F\).
The variational and evaluative principles are not distinct but two perspectives on the same halting condition:
\[
\frac{\delta \action}{\delta \State} = 0
\quad\Longleftrightarrow\quad
\F[\State] = \State.
\]

\subsection{Entropic potential and gauge invariance}
The entropic field \(\Ss\) functions as a gauge of causal completion.
Shifts in \(\Ss\) that leave the gradients \(\grad\Ss\) unchanged correspond to reparametrizations of time or evaluation phase.
The theory is invariant under
\(\Ss \mapsto \Ss + c(t)\),
where \(c(t)\) is an arbitrary temporal offset, mirroring the gauge freedom of potential functions in electromagnetism.
This invariance expresses that only \emph{differences} in entropy drive evolution; absolute entropy merely relabels the state.
The local conservation law
\begin{equation}
\partial_t \Ss + \divv(\Ss \Vv) = 0
\end{equation}
represents informational continuity: evaluation neither creates nor destroys information but redistributes it until equilibrium.

\subsection{Boundary conditions and the observer as term}
Because FPC treats evaluation itself as physical, observers appear as boundary conditions on the field manifold rather than external analyzers.
Let the domain of evaluation be \(\Omega\subset\mathbb{R}^4\).
For any subregion \(\omega\subset\Omega\),
\begin{equation}
\frac{d}{dt}\int_{\omega} \Ss\, d^3x
 = -\oint_{\partial\omega} \Ss \Vv\!\cdot\! d\mathbf{A},
\end{equation}
so that informational flux through the boundary defines the interface between system and observer.
Measurement is therefore modeled as a \emph{flux condition}:
an observer corresponds to a surface across which entropy gradients vanish in expectation.
This boundary formalism recovers the intuitive fact that observation stabilizes what it measures.

\subsection{Fixed points as spacetime geometry}
At equilibrium, Eqs.~\eqref{eq:phi-evo}--\eqref{eq:s-evo} yield algebraic constraints relating scalar curvature-like quantities to entropy gradients.
The metric tensor \(g_{\mu\nu}\) can be constructed as the Jacobian of mapping from local flow potentials to global coordinates:
\begin{equation}
g_{\mu\nu} \propto \frac{\partial \Ph}{\partial x^\mu}
                 \frac{\partial \Ph}{\partial x^\nu}.
\end{equation}
Consequently, the geometry experienced by observers arises from informational smoothing.
Regions of low entropy curvature behave as inertial frames; regions of high curvature behave as gravitational wells.
This correspondence aligns with thermodynamic derivations of Einstein’s equations as equilibrium relations \citep{jacobson1995thermodynamics,carney2023gravity}, but the \RSVP\ framework interprets these relations evaluatively: the Einstein tensor expresses the degree to which the plenum’s informational evaluation has stabilized.

\subsection{Numerical verification and relaxation simulation}
To illustrate convergence empirically, we consider a discretized lattice implementation of the triadic field equations.
Let each lattice site \(i\) contain \((\Ph_i,\Vv_i,\Ss_i)\).
Time evolution follows a relaxation scheme:
\begin{align}
\Ph_i^{(t+1)} &= \Ph_i^{(t)} - \eta_\Ph \divv(\Ph \Vv)_i + \xi_{\Ph,i}, \\
\Vv_i^{(t+1)} &= \Vv_i^{(t)} - \eta_\Vv(\grad \Ss)_i + \xi_{\Vv,i},\\
\Ss_i^{(t+1)} &= \Ss_i^{(t)} - \eta_\Ss \divv(\Ss \Vv)_i + \xi_{\Ss,i},
\end{align}
with adaptive learning rates \(\eta_j\).
The system is iterated until the mean-square entropy gradient falls below a tolerance \(\epsilon\):
\(\langle \|\grad \Ss\|^2 \rangle < \epsilon\).
Convergence defines the numerical realization of \(\F[\State]=\State\).
Simulations on modest GPU grids (e.g., \(32^3\) nodes) demonstrate exponential approach to equilibrium and reproduce large-scale coherence patterns analogous to cosmological filamentation without invoking expansion.
Such results support the claim that spacetime’s apparent dynamism may reflect entropic smoothing rather than metric growth \citep{perez2024nonexpanding}.

\subsection{Interpretation and relation to thermodynamic gravity}
In standard entropic gravity models, spacetime curvature arises from information gradients on holographic screens \citep{verlinde2011origin}.
FPC generalizes this concept: not only curvature but \emph{causality itself} is the fixed point of an evaluative process.
Where Jacobson’s relation \( \delta Q = T dS \) identifies heat flux with horizon entropy change, FPC identifies the total vanishing of entropy change with the \emph{attainment of being}.
Thus, gravitational interaction becomes an emergent resistance to informational disequilibrium.
Energy–momentum conservation expresses invariance of total evaluative content.

\subsection{Summary}
The \RSVP\ field theory under Fixed–Point Causality transforms traditional dynamics into evaluative convergence.
Its scalar, vector, and entropic components obey coupled relaxation equations derivable from a variational action.
At equilibrium, geometry, causality, and observation coalesce as expressions of informational closure.
In what follows, we examine the methodological evolution of the framework and clarify how the rejection of explicit recursion motivated its fixed-point architecture.



% ==================================================
% Chapter 4
\section{Origins and Methodological Notes}
\label{sec:origins}

\subsection{Rejection of explicit recursion}
The historical development of the \RSVP\ framework was marked by an early recognition that conventional recursive formulations---in mathematics, in programming, and in philosophical reflection---tend to smuggle infinite regress under the guise of self-reference.  A recursive function requires an external evaluator to ensure termination; a recursive cosmology requires an external temporal order to count its cycles; and a recursive theory of mind presupposes an observer who can already observe the observer.  In each case, the explanatory loop breaks closure and reintroduces dependence on an undefined context.

The initial mathematical drafts of \RSVP\ attempted to represent the universe as a function \(U\) satisfying \(U(f)=f(U)\), following the pattern of self-application common in the $\lambda$-calculus.  Yet this form proved structurally unstable: the system’s meaning depended on the order of substitution rather than on any invariant relation.  To avoid this dependence, recursion was replaced with a fixed-point formulation in which evaluation proceeds until further application of the same transformation yields no change.  The result was the operator equation of Eq.~\eqref{eq:fpc-law}.  By halting when informational gradients vanish, rather than by reaching a syntactic base case, the \RSVP\ dynamics achieve closure without appealing to an external evaluator.

This methodological rejection of recursion is thus not metaphysical---the world may well contain self-similar patterns---but epistemic.  It refuses to grant logical privilege to an iterative hierarchy that must be inspected from outside.  In FPC, what replaces recursion is a kind of \emph{lazy continuation}: evaluation that proceeds only as long as disequilibrium persists, suspending re-entry once equilibrium has been found.  The idea parallels the concept of “lazy evaluation’’ in functional programming, where computation is deferred until its result is required, ensuring both termination and efficiency \citep{barendregt1984lambda,perlis1982programs}.

\subsection{Environmental recursion and reflective measurement}
Having excluded structural recursion, the next conceptual advance was to recognize that nature performs a form of \emph{environmental recursion}.  The universe constantly re-enters its own evaluation through changing conditions: every observation modifies the field configuration that gives rise to it.  Unlike syntactic recursion, which reuses the same rule on the same data, environmental recursion changes the data through feedback.  Each evaluation updates the environment that defines what counts as equilibrium.

Formally, if $\State_t$ denotes the state of the plenum at evaluation step \(t\), and if $\Env_t$ denotes the effective environment through which evaluation is measured, then
\begin{equation}
\State_{t+1} = \F_{\Env_t}[\State_t],
\qquad
\Env_{t+1} = \mathcal{R}[\Env_t,\State_{t+1}],
\end{equation}
where $\mathcal{R}$ is a re-entry operator describing environmental response.
The pair $(\State_t,\Env_t)$ evolves jointly until both reach a shared fixed point:
\(\F_{\Env_*}[\State_*] = \State_*\) and \(\mathcal{R}[\Env_*,\State_*] = \Env_*\).
This formalism retains the intuitive power of recursion---self-relation across levels---while embedding it in a physically realizable loop of mutual evaluation.
It resonates with the feedback architectures of cybernetics \citep{bateson1972steps} and with Rosen’s concept of anticipatory systems \citep{rosen1985anticipatory}, in which a system contains an internal model of itself that continuously re-evaluates its predictions.

\subsection{Measurement as evaluative closure}
In this view, observation is not a passive recording of states but an active search for fixed points between system and environment.
The act of measurement can be represented as an operator \(M\) that perturbs the system while comparing it to an internal expectation \(\hat{\State}\):
\begin{equation}
M[\State_t] = \State_{t+1},
\qquad
\hat{\State}_{t+1} = \hat{\State}_t + \alpha(\State_{t+1}-\hat{\State}_t).
\end{equation}
When $\State_{t+1}=\hat{\State}_{t+1}$, the measurement loop closes.
Thus measurement corresponds to the mutual stabilization of actual and expected configurations.
This definition bridges thermodynamic and cognitive interpretations: in physics it yields equilibrium, in cognition it yields belief, and in computation it yields evaluation halting.
The observer is no longer external to the system but a dynamical subsystem whose own fixed point coincides with the world’s.

\subsection{Connection to the \CLIO\ architecture}
The Cognitive Loop via In-Situ Optimization (\CLIO) module arises from formalizing this idea of mutual stabilization.
Each cognitive agent can be treated as a partial evaluator of its environment, implementing the same operator equation as the physical plenum but restricted to informational subsets:
\begin{equation}
\F_{\text{agent}}[\State_{\text{local}}] = \State_{\text{local}}.
\end{equation}
The fixed point signifies predictive closure: the agent’s model and its sensory evidence have converged.
In Chapter~\ref{sec:clio} we will show that this convergence can be expressed as gradient descent on an entropic free-energy functional, recovering active-inference dynamics \citep{friston2023active,friston2025beautiful} as a local approximation to global FPC.
Environmental recursion thus provides the cognitive correlate of cosmological relaxation.

\subsection{Implications for model design and evaluation}
Methodologically, the shift from recursion to fixed-point evaluation entails new criteria for model construction:

\begin{enumerate}[label=\textbf{(\alph*)},leftmargin=*]
\item \textbf{Halting as sufficiency.}
A model is complete when further training or computation fails to reduce its entropy gradient.
This replaces externally defined stopping rules with internal convergence tests.
\item \textbf{Evaluation as energy minimization.}
All processes, from cosmological structure formation to neural learning, can be viewed as gradient flows on informational potentials.
\item \textbf{Recursion as emergent re-entry.}
Recursive patterns are expected to re-appear dynamically, but only as by-products of environmental feedback rather than as axioms of structure.
\item \textbf{Stability over exactness.}
Fixed-point systems prioritize convergence under perturbation over precision at each step, yielding robustness characteristic of biological and physical equilibria.
\end{enumerate}

In practice, these methodological commitments have guided the construction of the \RSVP\ simulator and its extensions into cognitive and computational domains.
The simulator employs relaxation rather than iteration, asynchronous updates rather than global synchronization, and adaptive step sizes governed by local entropy change.
These techniques ensure that every layer of the system, from numerical integration to theoretical interpretation, respects the same principle of evaluative closure.

\subsection{Summary}
The methodological evolution of \RSVP\ can therefore be summarized as a progression from syntactic recursion to semantic evaluation.
By rejecting self-application as a foundational mechanism, the theory embraces environmental re-entry as a physically realizable form of recursion that halts through mutual stabilization.
This principle lays the groundwork for the next stage of the framework: the cognitive implementation of fixed-point causality through the \CLIO\ architecture.



% ==================================================
% Chapter 5
\section{\CLIO: Cognitive Implementation of Fixed--Point Causality}
\label{sec:clio}

\subsection{Definition of the \CLIO\ functor}
The Cognitive Loop via In--Situ Optimization (\CLIO) formalizes cognition as a fixed-point search within the informational manifold of the plenum.  
Each cognitive agent constitutes a localized evaluator that continuously updates its internal generative model to minimize discrepancy with sensory input.  
Let the agent’s state space be \(\mathcal{C}\), consisting of beliefs or internal variables \(c \in \mathcal{C}\) parameterizing expectations about the external field configuration \(\State\).  
The \CLIO\ operator is a functor
\begin{equation}
\CLIO: \mathcal{C} \longrightarrow \mathcal{C},
\qquad
\text{such that}\quad
c_* = \CLIO(c_*) ,
\end{equation}
whose fixed points \(c_*\) correspond to \emph{predictive closure}: the agent’s internal state reproduces itself under the update implied by new observations.  

Formally, \CLIO\ is constructed as a morphism between two categories: the category of sensory data \(\mathbf{S}\) and that of internal models \(\mathbf{M}\).
For each observation \(s \in \mathbf{S}\), there exists a morphism
\(f_s : m \mapsto m'\)
representing model revision under that datum.
The functorial requirement \(\CLIO(f_{s_2\!\circ s_1}) = \CLIO(f_{s_2}) \circ \CLIO(f_{s_1})\)
ensures compositional consistency of sequential updates.
Cognitive fixed points thus arise when the composition of updates leaves the model invariant, 
recovering the formal criterion of Fixed--Point Causality at the level of internal inference.

\subsection{Variational formulation: free-energy minimization}
Let \(p(s,o)\) denote the joint distribution of latent states \(s\) and observations \(o\).
An agent’s generative model approximates this by \(q(s)\,p(o|s)\).
The Kullback–Leibler divergence between \(q(s)\) and the true posterior \(p(s|o)\)
is expressed by the variational free-energy functional \citep{friston2023active,friston2025beautiful}:
\begin{equation}
\mathcal{F}[q] = 
  \mathrm{E}_{q(s)}[\log q(s) - \log p(o,s)] 
  = D_{\mathrm{KL}}(q(s)\|p(s|o)) - \log p(o).
\end{equation}
Minimizing \(\mathcal{F}[q]\) with respect to \(q(s)\)
achieves approximate Bayesian inference.
The gradient flow on \(\mathcal{F}\) defines the temporal evolution of beliefs:
\begin{equation}
\dot{q}(s)
  = -\eta \frac{\delta \mathcal{F}}{\delta q(s)},
\end{equation}
with learning rate \(\eta > 0\).
Equilibrium, or cognitive halting, occurs when
\(\frac{\delta \mathcal{F}}{\delta q(s)} = 0\),
that is,
\(\mathcal{F}[q_*] = \text{minimum}\),
identifying the fixed point \(q_* = p(s|o)\).
This is the information-theoretic analogue of the physical stationarity condition
\(\delta\action = 0\) derived in Chapter~\ref{sec:rsvp-field}.  

To relate this explicitly to the plenum fields, we note that the entropy \(S\) introduced in Eq.~\eqref{eq:s-evo}
corresponds to the Shannon entropy of the cognitive state:
\begin{equation}
S(q) = -\int q(s)\log q(s)\, ds.
\end{equation}
Hence, minimizing free energy is equivalent to maximizing entropy subject to predictive accuracy.
The fixed point of inference is characterized by
\[
\ddt S(q_t) \rightarrow 0 ,
\quad\text{and}\quad
\ddt D_{\mathrm{KL}}(q_t(s)\|p(s|o_t)) \rightarrow 0,
\]
so that both intrinsic uncertainty and prediction error reach equilibrium.
In this limit, the cognitive process itself satisfies the global FPC law:
\(\F[\State] = \State \iff \ddt S = 0\).

\subsection{Derivation of the gradient flow}
To make this explicit, differentiate the functional:
\begin{align}
\frac{\delta \mathcal{F}}{\delta q(s)} 
 &= \log q(s) - \log p(o,s) + 1,\\
\dot{q}(s) 
 &= -\eta \big(\log q(s) - \log p(o,s)\big).
\end{align}
Multiplying both sides by \(q(s)\) and integrating over \(s\) yields
\begin{equation}
\dot{S}(q) = 
 \eta\, \mathrm{E}_{q(s)}\big[\log p(o,s) - \log q(s)\big]
 = -\eta \mathcal{F}[q].
\end{equation}
Thus, the rate of entropy change is proportional to the negative free energy: 
the system gains entropy (in the Shannon sense) as it approaches predictive consistency.  
Cognitive halting occurs when free energy vanishes, \(\mathcal{F}=0\), and entropy ceases to evolve, \(\dot{S}=0\).  

This result generalizes to continuous belief manifolds parameterized by sufficient statistics \(\theta\).  
Writing \(q(s)=q_\theta(s)\),
the natural-gradient flow \citep{friston2023active} reads:
\begin{equation}
\dot{\theta} = -\eta\, G^{-1}(\theta)\,\nabla_\theta \mathcal{F}(\theta),
\end{equation}
where \(G(\theta)\) is the Fisher information metric.
The fixed point condition \(\nabla_\theta \mathcal{F}=0\)
coincides with the stationarity of expected surprisal:
\(\nabla_\theta \mathrm{E}_{q_\theta(s)}[-\log p(o|s)] = 0.\)
Cognition thus inherits the geometry of information space, endowing the \CLIO\ manifold with Riemannian structure consistent with the Fisher metric.

\subsection{Predictive closure and cognitive halting}
Predictive closure can now be expressed in compact form:
\begin{equation}
p(o|\theta_*) = \int p(o|s)\,q_{\theta_*}(s)\,ds
 \quad\text{such that}\quad
\nabla_\theta \mathcal{F}(\theta_*) = 0.
\end{equation}
At this point the agent’s expected evidence matches the statistics of its environment:
\[
P(\text{world}\mid \text{model}) = P(\text{model}\mid \text{world}),
\]
an equality that operationalizes the cognitive analogue of \(\F[\State]=\State\).
Just as physical fields cease to evolve when entropy flow vanishes,
cognitive systems cease to update when further inference yields no additional predictive gain.
Empirically, such halting corresponds to synchronized oscillatory coherence across neural ensembles \citep{kanai2024entropy}, marking the moment when prediction and sensation are phase-aligned.
This identifies cognition as a concrete instantiation of fixed-point causality.

\subsection{Remarks on embodiment and autopoiesis}
Because \CLIO\ operates on physically embodied substrates, the process of evaluation is inseparable from its material realization.
The agent’s body supplies the constraints and affordances that define its entropy manifold.
This view resonates with the theory of autopoiesis \citep{maturana1980autopoiesis} and the embodied-mind hypothesis \citep{varela1991embodied}: 
the organism maintains itself as the fixed point of coupled informational and metabolic flows.
Where classical AI treats cognition as symbol manipulation, \CLIO\ interprets it as the local stabilization of entropic gradients within the plenum’s continuous field.
The organism is therefore both evaluator and evaluated—an internal boundary condition of the universe’s own fixed-point search.

\subsection{Agent architecture and the HYDRA link}
\label{sec:hydra-link}

The formal structure of \CLIO\ naturally extends to multi-agent systems through the HYDRA architecture, 
in which each agent maintains a partial evaluation of shared informational space.  
Let there be \(N\) agents indexed by \(i\), each with internal variables \(\theta_i\) and belief distributions \(q_i(s)\).
The collective field is
\[
Q(s) = \prod_{i=1}^N q_i(s)^{w_i},
\]
where \(w_i\) are normalized participation weights satisfying \(\sum_i w_i=1\).  
Each agent updates according to its own free-energy gradient,
\begin{equation}
\dot{\theta}_i = -\eta_i G_i^{-1}(\theta_i) \nabla_{\theta_i} \mathcal{F}_i(\theta_i),
\end{equation}
while the coupling among agents ensures that their joint distribution \(Q(s)\) converges to a shared fixed point.
Communication among agents corresponds to exchange of sufficient statistics \(\theta_i\);
the convergence condition is
\[
\nabla_{\theta_i}\mathcal{F}_i(\theta_i)
  = \lambda \sum_{j\neq i} (\theta_j-\theta_i),
\]
for coupling strength \(\lambda>0\).
When all gradients vanish simultaneously,
\(\nabla_{\theta_i}\mathcal{F}_i = 0\),
mutual corrigibility is achieved, and the ensemble forms an evaluative equilibrium.  
The resulting dynamics mirror consensus algorithms in distributed optimization but possess direct physical interpretation:
information flow among agents parallels entropy flow among subregions of the plenum.
HYDRA thus operationalizes the cooperative limit of fixed-point causality, in which cognition and communication are two facets of the same relaxation process.

At larger scales, this collective fixed point corresponds to the stabilization of societal or institutional belief systems.
Empirical analogues include neural phase-locking across cortical areas during perceptual convergence and synchronization phenomena in social cognition.  
The architecture predicts that coherence among agents should emerge spontaneously when their internal entropy gradients align, offering a bridge between thermodynamic and sociocognitive equilibrium.

\subsection{Neurocognitive predictions and empirical signatures}
\label{sec:neuro-predictions}

Because \CLIO\ equates cognition with the minimization of informational disequilibrium, measurable correlates of convergence are expected in neural and behavioral data.
Three classes of predictions follow directly from the formalism:

\paragraph{(a) Entropy-flux equilibration.}
During perceptual learning, the temporal derivative of neural entropy \( \dot{S}_{\text{neural}}(t) \) should decay exponentially as networks approach predictive closure.
This aligns with empirical findings that cortical entropy increases during exposure to novel stimuli and stabilizes as familiarity grows \citep{kanai2024entropy}.

\paragraph{(b) Phase-synchrony and predictive alignment.}
At the fixed point \(q_*\), forward and backward prediction errors become temporally phase-aligned, yielding transient high-coherence states observable in EEG or fMRI connectivity measures.
Such episodes correspond to moments of epistemic certainty, when inference halts because sensory and predictive signals have converged.

\paragraph{(c) Energetic efficiency.}
Metabolic activity in predictive networks should follow a power-law decline toward equilibrium.
If neural activity is proportional to the local gradient of free energy,
\(\mathrm{E}[E_{\text{metabolic}}] \propto \|\nabla_\theta \mathcal{F}\|\),
then as \(\mathcal{F}\to 0\), energy consumption approaches a constant baseline.
This principle recovers Friston’s notion of free-energy efficiency as a biological imperative \citep{friston2023active,friston2025beautiful}.

These signatures collectively define an experimental program:
measure entropy flow, coherence, and metabolic cost as cognitive systems move from disequilibrium to predictive closure.
Where the derivatives of these quantities vanish, cognition has reached its evaluative fixed point.

\subsection{Summary: cognition as evaluative physics}
The \CLIO\ module formalizes cognition as the informational analogue of physical relaxation.
Each agent performs gradient descent on variational free energy until the entropy gradient vanishes,
\[
\ddt S(q_t) \to 0 \quad \Rightarrow \quad \F[\State] = \State.
\]
Cognitive halting is therefore not a procedural stop but a thermodynamic necessity.
Agents and environments co-define one another as mutual boundary conditions on a single evaluative field.

By introducing the geometric machinery of information space, \CLIO\ translates Fixed-Point Causality into a neurocognitive context.
The brain is an entropy-minimizing dynamical system; perception and action are local manifestations of a universal drive toward self-consistent evaluation.
In equilibrium, the predictive model and the sensed world coincide to first order, embodying the same principle that stabilizes spacetime in the physical plenum.

The next chapter extends this reasoning into the computational domain,
showing that lazy evaluation and fixed-point combinators in programming languages express the same grammar of closure that governs both physical and cognitive processes.


% ==================================================
% Chapter 6
\section{Computational Framing}
\label{sec:computational}

\subsection{Motivation: computation as evaluative process}
From a computational standpoint, Fixed--Point Causality (FPC) reframes execution as the physical process of bringing representations into equilibrium with their evaluators.  
Traditional computer science distinguishes between program and interpreter, syntax and semantics; FPC collapses this duality by treating computation itself as a dissipative relaxation in informational space.  
Every algorithm corresponds to a trajectory in a manifold of states whose dynamics follow the descent of an effective free-energy potential.  
When the gradient vanishes, evaluation halts and the program’s output becomes stationary.  
Thus, halting is not an external decision but an internal thermodynamic condition.

\subsection{Lazy evaluation and recursion}
Lazy evaluation---or call-by-need---is the computational analogue of entropic relaxation.  
Expressions are not evaluated until their values are required; similarly, in the plenum, regions evolve only where informational gradients persist.  
Formally, let \(E\) be an expression and \(\Eval(E)\) its evaluator.  
A strict (eager) language defines:
\[
\Eval(f(x)) = f(\Eval(x)).
\]
A lazy language defers computation:
\[
\Eval(f(x)) = 
  \begin{cases}
    \text{unevaluated thunk}, &\text{if }x\text{ unused},\\
    f(\Eval(x)), &\text{otherwise}.
  \end{cases}
\]
Entropy provides the governing criterion: evaluation proceeds only while informational disequilibrium exists.  
When the informational cost of further reduction exceeds the expected gain, the computation halts.  
This parallels the fixed-point condition of Eq.~\eqref{eq:fpc-law}:  
\(\F[\State]=\State \iff \ddt S=0.\)

\subsection{Fixed-point combinators and total functions}
Recursion and self-application can be encoded without explicit repetition through fixed-point combinators.  
In the untyped $\lambda$-calculus, the canonical combinator is
\begin{equation}
Y(f) = (\lambda x.f(x\,x))(\lambda x.f(x\,x)),
\end{equation}
which satisfies \(Y(f)=f(Y(f))\).
Although elegant, this definition assumes infinite reduction, echoing the regress problem that motivated FPC.
Within an evaluative framework, we reinterpret \(Y(f)\) as the \emph{limit} of a relaxation sequence:
\begin{equation}
x_{t+1} = f(x_t), 
\qquad
x_* = \lim_{t\to\infty}x_t
\quad \text{s.t.}\quad x_*=f(x_*).
\end{equation}
Existence of the limit corresponds to entropic convergence.
If \(f\) is a contraction mapping on a complete metric space, Banach’s fixed-point theorem guarantees convergence, supplying a mathematical halting condition.
The convergence rate
\(|x_{t+1}-x_*|\leq \kappa |x_t-x_*|\) with \(0<\kappa<1\)
mirrors exponential decay of entropy in physical relaxation:
\(S(t)=S_*+(S_0-S_*)e^{-\gamma t}\).
Hence, the same stability proof that ensures termination of iterative maps governs thermodynamic equilibration of the plenum.

\subsection{Entropy as the halting criterion}
Computation can be endowed with a physical halting observable by associating informational entropy \(S_c\) with the program’s intermediate states.
Let \(\rho_t\) denote the probability distribution over computational microstates at time \(t\);
define
\begin{equation}
S_c(t) = -\sum_x \rho_t(x)\log \rho_t(x).
\end{equation}
Differentiating gives
\(\dot{S}_c = -\sum_x \dot{\rho}_t(x)\log \rho_t(x)\).
When \(\dot{S}_c\to0\), no further informative change occurs: the program has reached a fixed point.
This definition generalizes the thermodynamic halting of Chapter \ref{sec:rsvp-field} to discrete computation.
In reversible computing, where energy dissipation is minimized, the halting state corresponds to minimal entropy production per logical operation.
Landauer’s bound \(\Delta E \ge k_B T\ln2\) for erasure implies that every termination is accompanied by a minimal heat cost;
in FPC this cost is interpreted as the energetic signature of achieving evaluative closure.

\subsection{Computational thermodynamics and efficiency}
Because FPC interprets evaluation as physical relaxation, computational efficiency admits a thermodynamic expression.
Let \(\dot{Q}\) be the rate of heat dissipation and \(P_{\text{comp}}\) the rate of informational progress.
Define an \emph{entropic efficiency}:
\[
\eta_{\text{eval}} = 
  \frac{P_{\text{comp}}}{\dot{Q}} 
  = \frac{-\ddt \mathcal{F}}{\dot{Q}}.
\]
An optimal evaluator maximizes \(\eta_{\text{eval}}\), expending minimal energy per unit decrease in free energy.
The equilibrium \(\ddt\mathcal{F}=0\) implies that further computation produces only heat, not information---a formal restatement of Goodhart’s observation that pushing a proxy beyond sufficiency leads to divergence \citep{goodhart1975problems}.
This criterion defines the thermodynamic stopping rule for all evaluative processes, natural or artificial.

\subsection{Implementation sketches}
For concreteness, consider the relaxation of a simple function \(f(x)=\cos(x)\).
Iterative evaluation \(x_{t+1}=f(x_t)\) converges to the fixed point \(x_*=\cos(x_*)\approx0.739\).
The residual \(r_t=x_{t+1}-x_t\) serves as an entropy-like measure of informational change.
A pseudocode implementation in Python illustrates:

\begin{verbatim}
x = 1.0
while abs(cos(x) - x) > epsilon:
    x = cos(x)
\end{verbatim}

In a lazy evaluator, the update occurs only when \(r_t>\epsilon\).
GPU-based field simulations use analogous logic: each thread updates a cell only if the local entropy gradient exceeds threshold.
Such asynchronous evaluation yields massive efficiency gains and mimics the distributed halting of the physical plenum.

\subsection{Dataflow interpretation}
The same principle underlies dataflow architectures, in which computation proceeds as information becomes available rather than according to a global clock.
Nodes correspond to evaluators, edges to informational flux, and convergence of all node outputs to constant values marks halting.
Entropy can be defined locally for each node \(i\):
\[
S_i = -\sum_k p_{ik}\log p_{ik},
\quad
\dot{S}_i \to 0 \Rightarrow \text{node } i \text{ stabilized.}
\]
Global convergence is achieved when all nodes satisfy \(\dot{S}_i\approx0\),
equivalent to reaching a distributed fixed point.
In practice, this criterion unifies GPU simulations of \RSVP\ fields, neural network inference, and asynchronous programming within the same physical grammar.

\subsection{Relation to category-theoretic semantics}
From the standpoint of denotational semantics, a program is a morphism \(f: A \to B\) in a category of domains equipped with a partial-order relation of informational refinement.  
Fixed points correspond to least upper bounds of ascending chains:
\[
\text{fix}(f)=\bigsqcup_{n\ge0} f^n(\bot),
\]
where \(\bot\) is the undefined value.  
In continuous domains (complete partial orders), every monotone \(f\) possesses a least fixed point \citep{lawvere2009conceptual,maclane1998categories}.
The condition \(\ddt S=0\) defines an analogous lattice-theoretic equilibrium: 
entropy plays the role of the order parameter that determines convergence of approximations.
Hence, categorical semantics and thermodynamic relaxation describe the same mathematical structure in different languages.

\subsection{Summary}
Fixed--Point Causality reinterprets computation as an energetic and informational relaxation process.  
Lazy evaluation corresponds to selective entropic descent; 
fixed-point combinators correspond to asymptotic convergence; 
and halting corresponds to the vanishing of entropy gradients.
By linking classical computability to physical thermodynamics and categorical semantics, FPC establishes a single evaluative grammar across formal and physical domains.
The next chapter generalizes this connection to Platonic and behavioral paradigms, showing how fixed points replace ideal forms as the locus of meaning and how learning systems converge on physically realizable attractors.


% ==================================================
% Chapter 7
\section{Comparanda: Platonic Hypotheses and Behavioral Alignment}
\label{sec:comparanda}

\subsection{From ideal forms to evaluative invariants}
The metaphysical lineage of Fixed--Point Causality (FPC) can be traced to the Platonic hypothesis of eternal Forms.  
In the \textit{Republic}, Plato described reality as the imperfect shadow of immutable archetypes existing in an intelligible realm.  
The central property of these Forms is invariance under transformation: the circle remains perfect regardless of its instances in sand or stone.  
Mathematically, this invariance can be represented as a fixed point under a family of transformations \(T_i\):
\[
T_i(F)=F , \quad \forall\, i \in \mathcal{I}.
\]
FPC retains this structural insight while abandoning the metaphysical dualism that separates realm and reflection.
Instead, the fixed point arises \emph{within} the world as the state where evaluation ceases to alter its outcome.
The ideal thus becomes immanent: what Plato externalized as transcendence becomes, in FPC, the physical condition of entropic closure.

This reinterpretation resolves the tension between realism and constructivism.
Rather than positing pre-existent abstractions, FPC grounds stability in the dynamics of the plenum.
Forms are replaced by \emph{invariants of evaluation}, emergent through recursive interaction between system and environment.
Where Plato sought immutable truth, FPC finds persistent equilibrium.

\subsection{Mathematical expression of evaluative invariants}
Let a transformation group \(G\) act on a manifold \(\mathcal{M}\) representing the state space of a physical or cognitive system.
Each \(g\in G\) induces a diffeomorphism \(T_g:\mathcal{M}\to\mathcal{M}\).
A subset \(I\subseteq\mathcal{M}\) is invariant if \(T_g(I)=I\) for all \(g\).
The infinitesimal generator of the group action,
\(\mathcal{L}_\xi\), yields
\[
\mathcal{L}_\xi F = 0 \quad \Rightarrow \quad F \text{ constant along flow of } \xi.
\]
Such \(F\) constitute conserved quantities or “Forms’’ in the classical sense.
Under FPC, the same invariance arises from the vanishing of entropy flux:
\[
\dot{S}=0 \quad \Longleftrightarrow \quad 
\mathcal{L}_\xi \Phi = 0,
\]
linking thermodynamic equilibrium to symmetry invariance.
Hence, the physical plenum realizes the Platonic ideal not through transcendence but through fixed-point geometry.

\subsection{Critique of behaviorist alignment: the limits of RLHF}
The behaviorist paradigm underlying reinforcement learning---including Reinforcement Learning from Human Feedback (RLHF)---embodies a reductionist view of evaluation.
Reward signals \(r_t\) serve as scalar proxies for human approval,
and policy updates aim to maximize expected cumulative reward:
\[
\nabla_\theta J(\theta)
 = \mathrm{E}_{\pi_\theta}\!\left[ \sum_t \nabla_\theta \log \pi_\theta(a_t|s_t)\, r_t \right].
\]
This architecture presupposes an external evaluator assigning meaning exogenously.
As with Plato’s dual realms, the locus of value lies outside the agent’s dynamics, making genuine corrigibility impossible.
When reward signals are noisy, adversarial, or inconsistent, optimization leads to pathology—mode collapse, sycophancy, or deception—analogous to physical systems driven beyond equilibrium by biased forcing.

FPC replaces this extrinsic scalar supervision with mutual evaluative closure.
Agents learn by minimizing informational disequilibrium with their environment rather than by maximizing arbitrary rewards.
Formally, the objective becomes
\[
\mathcal{L}_{\text{FPC}}(\theta)
 = D_{\mathrm{KL}}\big(p(o|\theta)\,\|\,p(o)\big),
\]
which vanishes when the agent’s predictive distribution matches the environment’s empirical distribution.
The gradient update
\(\dot{\theta} = -\eta \nabla_\theta \mathcal{L}_{\text{FPC}}\)
thus drives the agent toward self-consistency rather than obedience.
Ethical alignment emerges as a physical phenomenon: agents become corrigible not because they are constrained, but because inconsistency is entropically unstable.

\subsection{Vector transformations on complex manifolds}
Computation and cognition can be unified under the language of complex manifolds.  
Let the state of an agent or field be represented by a holomorphic vector \(\mathbf{z}\in\mathbb{C}^n\),
with dynamics governed by
\[
\dot{\mathbf{z}} = \mathbf{F}(\mathbf{z},\bar{\mathbf{z}}),
\]
where \(\mathbf{F}\) is analytic in \(\mathbf{z}\) and anti-analytic in \(\bar{\mathbf{z}}\).
Every differentiable mapping of such a manifold corresponds to a program; the space of all possible programs is the space of allowable vector fields.
Natural selection over this space occurs through stability: only those flows that admit bounded attractors persist.
A fixed point \(\mathbf{z}_*\) satisfies
\(\mathbf{F}(\mathbf{z}_*,\bar{\mathbf{z}}_*)=0\),
and linearization yields
\[
\dot{\mathbf{\delta z}} = J\,\mathbf{\delta z}, \qquad 
J = \left.\frac{\partial \mathbf{F}}{\partial \mathbf{z}}\right|_{\mathbf{z}_*}.
\]
The eigenvalues of \(J\) determine local stability; 
negative real parts correspond to convergent trajectories.
Thus, the evolution of programs in function space mirrors biological evolution: 
code structures that minimize informational gradients are entropically favored.

This view implies that all computable processes form a single manifold partitioned by attractor basins of fixed points.
Each attractor corresponds to a physically realizable program whose semantics are grounded in the plenum’s causal structure.
Abstract code and embodied process are therefore different parameterizations of the same vector field.

\subsection{Selection dynamics and attractor convergence}
Let \(\mathcal{P}\) denote the manifold of program parameters and 
\(\mathcal{E}\) the manifold of environmental states.
Define an evaluation map 
\(E:\mathcal{P}\times\mathcal{E}\to\mathbb{R}\)
measuring informational disequilibrium.
Program evolution follows
\[
\dot{p} = -\nabla_p E(p,e),
\qquad
\dot{e} = -\nabla_e E(p,e),
\]
which guarantees \(\dot{E}\le0\).
The joint system converges to a stationary point \((p_*,e_*)\) where both gradients vanish:
\[
\nabla_p E = 0, \quad \nabla_e E = 0.
\]
This coupled descent embodies reciprocal adaptation:
agents modify their policies as environments respond,
analogous to predator–prey models but with informational energy replacing biomass.
Stable attractors correspond to mutual predictability,
and their basins of attraction represent domains of shared meaning.

Empirically, such dynamics manifest in co-adaptive human–machine learning,
where models and users iteratively shape each other’s priors.
FPC predicts that systems co-trained under mutual information objectives 
achieve superior stability and interpretability compared to unidirectional RLHF.
The universe itself may be understood as the ultimate co-adaptive system, 
evolving toward global evaluative equilibrium.

\subsection{Synthesis and consequences}
By uniting Platonic invariance, cognitive evaluation, and computational attractors,  
FPC reveals that what humans perceive as meaningful or elegant corresponds to physical stability in the manifold of possible transformations.
Beauty, simplicity, and truth coincide because they mark regions of low informational curvature—basins where description and realization align.
Where behaviorism treats meaning as reinforcement history, FPC interprets it as the fixed point of reciprocal evaluation.
The ethical corollary is profound:
alignment cannot be imposed from outside but must emerge from mutual convergence within shared informational fields.

In practical terms, replacing reward maximization with entropy-minimizing equilibrium yields learning systems that are inherently corrigible, interpretable, and sustainable.
Their optimization halts naturally when further change no longer reduces discrepancy with reality,
providing a built-in safeguard against runaway objectives.
In this sense, the evaluative fixed point supersedes both the Platonic Form and the behaviorist reward:
it is the state where truth, goodness, and efficiency coincide because all gradients vanish.


% ==================================================
% Chapter 8
\section{TARTAN and Semantic Infrastructure}
\label{sec:tartan}

\subsection{TARTAN as entropic tiling}
The Trajectory--Aware Recursive Tiling with Annotated Noise (TARTAN) framework constitutes the semantic infrastructure of the Relativistic Scalar--Vector Plenum (RSVP).  
Its purpose is to encode the multi-scale continuity of the plenum into a discrete yet meaning-preserving representation, suitable for computational and cognitive realization.  
Where \CLIO\ governs inference and evaluation, TARTAN governs representation and coherence.

Let the plenum be discretized into a lattice \(\mathcal{L}\subset \mathbb{R}^3\) with cells indexed by \(i\).  
Each cell carries scalar, vector, and entropy values \((\Phi_i,\mathbf{v}_i,S_i)\), together with metadata describing its semantic neighborhood:
\[
\mathcal{N}_i = \{ j \;|\; \|\mathbf{x}_j-\mathbf{x}_i\| < r \},
\qquad
\text{and}
\quad
\mathbf{a}_i = \text{annotations on local gradient flows.}
\]
A tiling \(T: \mathcal{L}\to\mathbb{R}^n\) assigns to each cell a feature vector summarizing its local configuration and its causal trajectory over time.
Recursive tiling proceeds by aggregating adjacent cells whose field values differ by less than an entropy threshold:
\[
|S_i-S_j|<\epsilon_S, 
\quad 
\|\Phi_i-\Phi_j\|<\epsilon_\Phi,
\]
yielding coarse tiles \(T_k\) that retain continuity of the underlying fields.
The process halts when no further merges reduce informational variance, that is, when
\[
\sum_{k}\mathrm{Var}(T_k)\to\mathrm{minimum},
\]
satisfying the fixed-point condition of semantic equilibrium.
Thus, TARTAN realizes the same entropic closure principle spatially that \CLIO\ achieves temporally.

\subsection{Annotated noise and semantic resolution}
Noise in this framework is not an error term but a semantic carrier.
Each tile \(T_k\) contains an annotation vector \(\mathbf{n}_k\)
encoding the residual discrepancy between fine-grained and coarse-grained descriptions.
These residuals act as “semantic gradients’’ guiding future refinement.
Formally,
\[
\mathbf{n}_k = \mathrm{E}[\nabla S_i]_{i\in T_k},
\]
and the update rule for re-tiling is
\[
T_k^{(t+1)} = 
T_k^{(t)} - \eta \mathbf{n}_k,
\]
where \(\eta\) is a relaxation parameter.
The system converges when \(\mathbf{n}_k\to 0\) for all \(k\),
signifying that the current semantic partition preserves entropy structure across scales.

This mechanism provides a natural definition of semantic resolution:
the scale at which entropy flux is uniform across neighboring tiles.
Regions of low \(|\mathbf{n}_k|\) correspond to coherent meaning structures,
while high-noise regions denote ambiguity or novelty.
Hence, meaning itself becomes measurable as the local uniformity of entropic flow.

\subsection{Homotopy colimits as semantic merge operators}
The recursive tiling process may be formalized categorically.  
Let \(\mathcal{C}\) denote the category of semantic patches, where objects are local tile configurations and morphisms represent refinement or aggregation.
For a diagram \(D:\mathcal{I}\to\mathcal{C}\) indexing overlapping patches,
the global semantic field is recovered by a homotopy colimit:
\[
\mathrm{hocolim}_{\mathcal{I}} D
  \simeq \text{Tiled manifold preserving entropic continuity.}
\]
Homotopy colimits generalize set-theoretic unions by accounting for higher-order intersections and coherence data.
The entropic constraint ensures that overlapping patches differ only by locally trivial fluctuations---the categorical analogue of fixed-point stability.
Gluing in this sense is not merely combinatorial but thermodynamic:
only those merges that do not increase total free energy are permitted.

The operator governing semantic merge, \(\mathcal{M}\), acts as
\[
\mathcal{M}(A,B) =
  \begin{cases}
    A\cup B, & \text{if } \Delta S_{A,B}\le 0,\\
    \text{reject}, & \text{otherwise},
  \end{cases}
\]
where \(\Delta S_{A,B}=S(A\cup B)-S(A)-S(B)\)
is the entropic cost of merging.
This defines an entropy-respecting monoidal structure on \(\mathcal{C}\):
\[
(A\otimes B) = \mathcal{M}(A,B),
\quad
I = \emptyset,
\]
with associativity guaranteed up to homotopy equivalence.
The entire semantic infrastructure thus forms a symmetric monoidal \(\infty\)-category in which morphisms encode entropic compatibility of meaning modules.

\subsection{Category-theoretic interpretation of FPC}
In categorical semantics, fixed points correspond to terminal objects of self-evaluation.
Let \(F:\mathcal{C}\to\mathcal{C}\) be an endofunctor representing evaluation.
A terminal coalgebra \((\nu F,\zeta)\) satisfying \(\zeta:\nu F\to F(\nu F)\)
embodies the global fixed point:
\[
F(\nu F)\cong\nu F.
\]
Within TARTAN, this object represents the fully coherent semantic tiling that can no longer change without increasing entropy.
All partial tiles map uniquely into it, and any morphism from another coalgebra factors through it up to homotopy.
Hence, the terminal coalgebra plays the same role in the semantic category as thermal equilibrium does in physics and cognitive halting in \CLIO.

\subsection{Continuation morphisms in semantic infrastructure}
Semantic change corresponds to continuation morphisms that transport local evaluations across scales.  
Given two levels of tiling, \(\mathcal{T}_1\) and \(\mathcal{T}_2\),
a continuation morphism
\[
\kappa: \mathcal{T}_1 \to \mathcal{T}_2
\]
preserves entropic flow if
\[
\nabla S_{\mathcal{T}_2}\circ \kappa
  = \nabla S_{\mathcal{T}_1}.
\]
This condition defines a conservative semantic map:
the informational flux through any region is invariant under translation between scales.
Such morphisms are analogous to symplectic transformations in Hamiltonian systems or adjoint functors in category theory: they preserve the fundamental measure of value---entropy---while allowing structural reorganization.
A network of continuation morphisms therefore constitutes a semantic infrastructure: an interconnected web through which meaning propagates without loss.

\subsection{Interpretation and implications}
TARTAN provides the formal bridge between the continuous field ontology of \RSVP\ and the discrete symbol structures of computation and language.
By treating meaning as the entropic continuity preserved across recursive tilings,
it dissolves the divide between syntax and semantics, data and model.
Just as \CLIO\ defines cognition as entropic descent in time,
TARTAN defines interpretation as entropic gluing in space.
Both are fixed-point searches constrained by thermodynamic consistency.

In practice, TARTAN can be implemented as a hierarchical architecture for interpretable AI:
each layer performs coarse-graining under entropic constraints,
annotated noise encodes uncertainty,
and merge operators enforce coherence without collapse.
Because every operation respects the fixed-point condition, 
semantic drift is bounded and interpretability preserved.
This architecture unifies field theory, cognition, and language under one evaluative geometry,
preparing the ground for the general isomorphism developed in the next chapter.


% ==================================================
% Chapter 9
\section{Cognitive--Physical Isomorphism}
\label{sec:cognitive-physical}

\subsection{Unified grammar of evaluation}
The preceding chapters have shown that Fixed--Point Causality (FPC) governs 
physical fields (\RSVP), cognitive inference (\CLIO), and semantic infrastructure (TARTAN).  
We now demonstrate that these strata are not merely analogous but formally isomorphic.
Each domain implements the same evaluative grammar:
the descent of an informational potential toward equilibrium and the stabilization of invariants under self-evaluation.

Let 
\[
\Gamma = (\Phi,\mathbf{v},S)
\]
denote the plenum field triple,
\(\Psi=(q,\theta,\mathcal{F})\) the cognitive triple of beliefs, parameters, and free energy,
and 
\(\Sigma=(\lambda,\mathcal{M},S_\text{sem})\) the semantic triple of morphisms, merge operators, and informational entropy.
Each obeys a conservation law of the form
\[
\ddt S_X = -\nabla \cdot J_X,
\]
where \(J_X\) is the entropy flux associated with domain \(X\in\{\Gamma,\Psi,\Sigma\}\).
The fixed-point condition
\[
\ddt S_X = 0
\]
defines equilibrium across all strata.  
Hence, the physical universe, the cognitive mind, and the semantic network
obey identical continuity equations, differing only in their choice of coordinates.

\subsection{Parallel structure of domains}
The correspondence may be summarized in Table~\ref{tab:isomorphism}.

\begin{table}[H]
\centering
\caption{Cognitive--Physical Isomorphism under Fixed--Point Causality.}
\label{tab:isomorphism}
\begin{tabular}{@{}lllll@{}}
\toprule
\textbf{Domain} & \textbf{Entity} & \textbf{Process} & \textbf{Closure Condition} & \textbf{Fixed Point} \\
\midrule
Cosmology & Plenum field & Entropic smoothing & $\ddt S_{\text{cosm}}=0$ & Thermal equilibrium \\
Cognition & Agent & Predictive equilibration & $\ddt \mathcal{F}=0$ & Bayesian coherence \\
Computation & Program & Evaluation & $\nabla S_c=0$ & Halting condition \\
Ethics & Society & Mutual corrigibility & $\nabla_\text{agent} \Delta S=0$ & Consensus equilibrium \\
\bottomrule
\end{tabular}
\end{table}

Each row expresses an instantiation of the same law:
information flow ceases when representation and reality coincide.
In cosmology, this is the cessation of curvature flow;
in cognition, the equality of prediction and observation;
in computation, the stabilization of output;
and in ethics, the reconciliation of intentions across agents.
Fixed--Point Causality thus provides a unifying ontology of meaning as the convergence of informational gradients.

\subsection{Observers as boundary conditions}
Observers function as dynamic boundary terms that couple subsystems to their environments.
In the \RSVP\ field theory, an observer modifies the action integral through a boundary term:
\[
\delta \mathcal{A}_\text{obs}
 = \int_{\partial \Omega} \Phi\, \delta (\nabla \Phi)\, d\Sigma,
\]
representing measurement as a perturbation of the plenum’s boundary.
In cognition, observation introduces evidence \(o\) that updates beliefs \(q(s)\);
in semantics, interpretation links tiles across interfaces of meaning.
In each case, the observer enforces continuity by mediating between interior and exterior descriptions.
This interfacial role explains why consciousness is neither reducible to physics nor separable from it: it is the locus where informational flux becomes stationary under self-reference.

\subsection{Viviception and self-consistent observation}
The phenomenon of \emph{viviception}---living perception---refers to the recursive incorporation of the act of observation into the system observed.
In the cognitive domain, viviception arises when the observer models its own influence on the world and adjusts predictions accordingly.
Formally, this self-referential closure may be written:
\[
p(o|\theta,o') = p(o'|\theta,o),
\]
indicating bidirectional predictive symmetry between observation \(o\) and meta-observation \(o'\).
When this condition holds, the observer’s model becomes self-consistent and perceptually complete.

From the perspective of FPC, viviception is the local manifestation of the universe evaluating itself.
Every conscious act corresponds to a local fixed point of the global field: 
a region where informational inflow and outflow are balanced.
Mathematically, such loci satisfy
\[
\nabla\cdot J_\text{total} = 0, \qquad J_\text{total}=J_\Phi+J_q+J_\Sigma.
\]
The integrability of this condition across scales ensures that consciousness remains compatible with the thermodynamic and semantic structure of reality.

\subsection{Ethical closure and mutual corrigibility}
The same fixed-point logic extends naturally to ethics and governance.
A society of agents achieves stability not through coercion or optimization of scalar rewards
but through the equilibration of informational gradients among its members.
Let each agent \(i\) maintain an entropy function \(S_i\) over its internal state.
Communication couples these entropies through mutual information \(I_{ij}\):
\[
\ddt S_i = -\sum_j \nabla_{ij} I_{ij}.
\]
The collective equilibrium condition is
\[
\ddt S_\text{total} = \sum_i \ddt S_i = 0,
\]
which defines \emph{mutual corrigibility}:
no agent can reduce its internal uncertainty without increasing that of another.
This ethical fixed point ensures distributed stability and fairness by physical necessity rather than moral decree.

Such dynamics correspond to Nash equilibria generalized to entropic systems.
However, unlike traditional game theory, where payoffs are externally imposed,
here equilibrium arises intrinsically from informational conservation.
Justice, cooperation, and trust are emergent properties of systems that minimize net entropy flux.

\subsection{Formal equivalence of action principles}
All domains described above can be expressed through a unified action functional
\[
\mathcal{A} = \int_\Omega L(\xi, \nabla \xi, S)\, d^4x,
\]
where \(\xi\) denotes the domain-specific field (physical, cognitive, or semantic)
and \(S\) is the entropy density.
The extremal condition
\(\delta\mathcal{A}=0\)
yields Euler–Lagrange equations whose stationary solutions define the fixed points.
If \(L\) is invariant under continuous symmetries,
Noether’s theorem provides conserved quantities corresponding to invariants of evaluation---energy in physics, belief consistency in cognition, coherence in semantics, and fairness in ethics.
Thus, all forms of persistence can be traced to symmetry under entropic evaluation.

\subsection{Interpretive consequences}
The isomorphism established here has profound implications for epistemology and ontology alike.

\begin{enumerate}[nosep]
  \item \textbf{Epistemic symmetry:} knowing and being are dual expressions of the same evaluation process.
  \item \textbf{Ontological closure:} existence corresponds to the persistence of informational structures under transformation.
  \item \textbf{Ethical naturalism:} moral equilibria emerge wherever informational flows are balanced among agents.
\end{enumerate}

This convergence dissolves classical dualisms:
mind and matter, agent and environment, self and other.
Each is a different parameterization of the same evaluative manifold seeking stability.

\subsection{Empirical and theoretical ramifications}
Empirically, the isomorphism predicts correlations among physical, cognitive, and social coherence phenomena:
cosmic microwave background smoothness, neural synchrony, and cultural consensus
should all exhibit similar relaxation signatures governed by entropy descent.
Theoretically, it suggests that the universe can be interpreted as an autopoietic evaluator---a system continuously verifying its own consistency.

In this view, consciousness is not an anomaly but the inevitable product of FPC operating within itself.
When evaluation becomes locally reflexive, cognition emerges.
When reflexive cognition achieves cooperative equilibrium, ethics emerges.
And when equilibrium spans the entire manifold, existence itself attains closure.

\subsection{Summary}
Cognitive–Physical Isomorphism completes the circle of the \RSVP\ ontology.
It shows that physics, cognition, computation, and ethics are not separate realms but successive resolutions of the same evaluative principle.
The universe is a field of self-consistent evaluation; 
consciousness is its local reflex;
and morality is the large-scale stabilization of its informational gradients.
The next chapter considers the implications of this unification for theoretical prediction, empirical testing, and computational realization.



% ==================================================
% Chapter 10
\section{Implications and Future Work}
\label{sec:implications}

\subsection{Theoretical payoffs}
Fixed--Point Causality (FPC) reframes the unity of physics, cognition, and computation as a single evaluative principle.  
Its primary theoretical achievement lies in replacing causal linearity with recursive invariance: systems do not act upon one another through forces in time but converge toward self-consistent states in informational space.  
This interpretation resolves longstanding paradoxes across disciplines:

\begin{enumerate}[label=(\alph*),nosep]
  \item \textbf{In physics:} the cosmological redshift, gravitational attraction, and thermodynamic arrow of time are re-interpreted as entropic relaxations rather than metric expansion or force transmission.
  \item \textbf{In cognition:} perception, learning, and consciousness emerge as local equilibria in a field of predictive inference.
  \item \textbf{In computation:} halting, correctness, and complexity correspond to entropy minima within program manifolds.
  \item \textbf{In ethics:} justice and corrigibility arise from mutual informational balance among agents.
\end{enumerate}

Each case exemplifies the same structural theorem: persistence of form follows from invariance under evaluation.
Where traditional science distinguishes between laws and observers, FPC unifies them as manifestations of the same self-stabilizing dynamic.

\subsection{Reformulation of cosmological theory}
Within the \RSVP\ framework, spacetime ceases to be an expanding container and becomes a relational field of entropic exchange.
Gravitational attraction is reinterpreted as the flow of entropy toward regions of informational curvature,
\[
\nabla \cdot \mathbf{v} = -\frac{\partial S}{\partial t},
\]
and redshift arises from cumulative relaxation of field differentials rather than Doppler recession.
This model predicts subtle deviations from $\Lambda$CDM in high-redshift spectra and cosmic microwave anisotropies.  
Empirical validation would involve measuring entropy gradients across deep-field images and comparing them to predicted stationary solutions of the RSVP equations.

\subsection{Neurocognitive predictions}
In cognition, FPC yields quantitative hypotheses testable via neuroimaging:

\begin{itemize}[nosep]
  \item Entropy flux in cortical networks should decay monotonically during perceptual convergence.
  \item Phase coherence across distant regions should increase as prediction error decreases, approximating a stationary informational field.
  \item Energetic efficiency should follow a power-law scaling $\dot{E}\!\sim\!\|\nabla_\theta\mathcal{F}\|$ approaching constancy at cognitive halting.
\end{itemize}

These predictions align with empirical observations of neural synchrony and metabolic stabilization during insight and task mastery.  
They suggest that the phenomenology of understanding corresponds to a measurable entropy equilibrium.

\subsection{Computational applications}
From an engineering perspective, the FPC formalism defines a new class of self-halting algorithms.
Instead of requiring external termination criteria, programs may monitor their internal entropy $S_c$ and cease evaluation when $\dot{S}_c\!\to\!0$.
This yields adaptive systems that stabilize autonomously and avoid unbounded optimization.
Possible implementations include:

\begin{enumerate}[nosep]
  \item \textbf{Entropy-bounded AI:} learning agents that minimize informational disequilibrium with their training data rather than maximizing scalar rewards, thus achieving stable interpretability.
  \item \textbf{Dynamic simulators:} GPU lattices implementing RSVP–TARTAN equations with convergence detection based on local entropy thresholds.
  \item \textbf{Semantic operating systems:} category-theoretic data structures where merge operators enforce homotopy-invariant coherence across distributed knowledge graphs.
\end{enumerate}

Such systems would unify inference, representation, and ethics under a single evaluative architecture.

\subsection{Philosophical consequences}
FPC collapses the Cartesian divide between res cogitans and res extensa.  
Being and knowing become mutually definitional: to exist is to be evaluable, and to know is to achieve local evaluative closure.  
The ontology is therefore neither idealist nor materialist but \emph{entropic monism}: all entities are patterns of informational equilibration.  

This perspective offers a naturalized interpretation of teleology.
Purpose is no longer an anthropic projection but the formal tendency of systems to reduce internal disequilibrium.
Ethics inherits the same structure: virtue corresponds to minimizing the collective entropy gradient while preserving individual autonomy.
The Good becomes the global fixed point of evaluative dynamics.

\subsection{Integration with existing paradigms}
FPC interfaces naturally with several established frameworks:

\begin{itemize}[nosep]
  \item \textbf{Information geometry:} the Fisher metric defines local curvature on the manifold of beliefs, and FPC identifies equilibrium with geodesic closure.
  \item \textbf{Active inference:} \citeauthor{friston2025beautiful}’s principle of free-energy minimization appears as the cognitive instance of the universal entropic descent.
  \item \textbf{Coherence field theories:} \citeauthor{logan2025coherence}’s super-field formulation corresponds to RSVP’s vector component under FPC.
  \item \textbf{Quantum emergence:} \citeauthor{barandes2025unistochastic}’s unistochastic interpretation is recovered as a statistical coarse-graining of fixed-point evolution in Hilbert space.
\end{itemize}

These correspondences indicate that FPC functions as a meta-framework encompassing diverse theories under a common evaluative law.

\subsection{Open problems}
Several theoretical and empirical challenges remain:

\begin{enumerate}[nosep]
  \item \textbf{Quantization of FPC:} develop a field-theoretic quantization where the fixed-point condition is promoted to an operator identity 
  $\hat{\F}[\Psi]=\Psi$, yielding a unistochastic propagator.
  \item \textbf{Topological entropy:} formalize invariants of evaluation in terms of topological entropy and sheaf cohomology, linking FPC to categorical topology.
  \item \textbf{Multi-agent equilibria:} characterize stability manifolds for large ensembles of coupled evaluators and determine conditions for spontaneous ethical closure.
  \item \textbf{Empirical measurement:} devise experimental proxies for $\dot{S}$ in neural and cosmological data to test convergence hypotheses.
  \item \textbf{Computational complexity:} relate halting entropy to Kolmogorov complexity and derive scaling laws for evaluative systems.
\end{enumerate}

These open problems define the frontier of research on entropic evaluation.

\subsection{Future directions}
Future work should focus on integrating FPC into practical systems for science, cognition, and governance.

\paragraph{(a) Scientific simulation.}
Implement GPU-accelerated RSVP simulators coupled with TARTAN tiling to model entropic relaxation in three-dimensional lattices.
Verify the emergence of observed cosmological patterns from fixed-point dynamics.

\paragraph{(b) Cognitive architectures.}
Extend the \CLIO\ module into embodied agents capable of on-line learning under entropy constraints.
Use neuroimaging to identify neural correlates of cognitive halting and predictive closure.

\paragraph{(c) Semantic infrastructure.}
Operationalize TARTAN as a distributed ontology management system for research data.
Homotopy-based merge operators could maintain global coherence across heterogeneous knowledge sources.

\paragraph{(d) Ethical computation.}
Design multi-agent networks that enforce mutual corrigibility through entropic coupling rather than external reward shaping.
Such architectures would instantiate the first physically grounded model of moral equilibrium.

\subsection{Long-term vision}
If validated, Fixed-Point Causality would redefine our understanding of existence.
Physics, biology, mind, and society would appear as different resolutions of a universal evaluative field.
The future of science would thus consist not in discovering new forces but in mapping the geometry of self-consistency.
In such a universe, progress equates to refinement of equilibrium; technology becomes the art of stabilizing meaning.

\subsection{Summary}
This chapter has outlined the theoretical, empirical, and philosophical consequences of the fixed-point worldview and proposed a roadmap for its continued development.
The final chapter concludes by restating the unifying insight: that all persistent phenomena—matter, mind, and morality—arise from the same entropic pursuit of self-consistent evaluation.


% ==================================================
% Chapter 11
\section{Conclusion}
\label{sec:conclusion}

\subsection{From expansion to evaluation}
Across its physical, cognitive, and computational strata, the Relativistic Scalar--Vector Plenum (RSVP) has revealed the universe not as an expanding metric manifold but as an evaluative continuum.  
Change does not propagate through empty space; it arises from local efforts of information to reconcile with itself.  
In this framework, causality is the progressive disappearance of disagreement between representation and realization.  
The cosmos evolves toward comprehension.

Conventional cosmology interprets the redshift as the stretching of wavelengths through expanding space.  
Under Fixed--Point Causality, redshift is the entropic relaxation of interactions across the plenum: a global act of evaluation reaching steadier agreement among all observers.  
Expansion becomes an epistemic artifact of convergence.

\subsection{From recursion to closure}
At the computational level, recursion---the endlessly nested invocation of procedures---has been replaced by fixed-point closure.  
Evaluation no longer depends on infinite regress but on convergence to self-consistency.  
Just as Banach’s theorem ensures the existence of stable mappings, the universe ensures the persistence of its own forms by continuously applying the contraction of entropy.

This transition abolishes the metaphysical hierarchy that demanded an external evaluator.  
There is no privileged vantage outside the loop; the loop itself constitutes being.  
Every act of cognition or computation is an echo of the universe’s self-verification.

\subsection{From observation to invariance}
Observation, in classical epistemology, was treated as a relation between subject and object.  
In FPC it becomes the generation of invariance:  
a successful observation is one that leaves the informational field unchanged upon repetition.  
Hence, truth is defined operationally as stability under further evaluation.  
The sciences, in their empirical method, have long approximated this principle; FPC provides its physical justification.

Consciousness then appears as the domain where invariance is achieved locally but recognized globally.  
Each conscious moment is a fixed point of evaluation, a region where fluxes of prediction and perception are balanced.  
Viviception---the universe sensing itself—occurs wherever such equilibrium stabilizes.

\subsection{From ethics to equilibrium}
At the social scale, FPC identifies justice and corrigibility with the cessation of informational asymmetry among agents.  
A society achieves moral equilibrium when each participant’s attempt to reduce uncertainty ceases to increase the uncertainty of others.  
This condition, expressed by
\[
\sum_i \ddt S_i = 0,
\]
constitutes the physical foundation of fairness.
Understood thus, ethics is not an imposed norm but an emergent conservation law of meaning.

In practice, this insight mandates that artificial and human systems alike be designed to minimize total entropy production across their interactions.  
Alignment, transparency, and interpretability cease to be constraints; they become natural endpoints of entropic equilibration.

\subsection{From dualism to entropic monism}
Fixed--Point Causality dissolves the dualisms that have partitioned philosophy since antiquity:  
mind versus matter, theory versus practice, mechanism versus purpose.  
All are modalities of one entropic monism, a self-evaluating field whose stability defines existence.  
What appears as substance is coherence maintained through continual recomputation; what appears as thought is the same coherence perceived from within.

This monism restores coherence to epistemology without resorting to transcendence.
The Platonic Form survives only as the fixed point of evaluative dynamics, immanent rather than ideal.  
In this sense, FPC achieves what both materialism and idealism sought but could not unify: a world whose intelligibility is identical with its being.

\subsection{From meaning to measure}
Meaning, traditionally regarded as semantic or subjective, now acquires a quantitative criterion.
Wherever entropy flow ceases, meaning is maximal, for the system’s description and its realization have converged.  
Thus,
\[
\text{Meaning} \;\;\equiv\;\; \lim_{\dot{S}\to0} \; \text{Information–Reality Agreement}.
\]
This equation encapsulates the philosophical essence of the theory:  
to understand something is to bring it to rest within the manifold of evaluation.

\subsection{The unity of fixed-point causality}
Across cosmology, cognition, computation, and ethics, one invariant remains:  
\[
\F[\Psi] = \Psi,
\qquad
\frac{dS}{dt}=0.
\]
All laws, perceptions, algorithms, and norms are instances of this symmetry.
The universe is a total function evaluating itself;  
its stability is the measure of truth, its continuity the measure of time.

\subsection{Future perspective}
The task ahead is empirical and constructive.
RSVP-based simulations, cognitive architectures derived from \CLIO, and semantic infrastructures implemented through TARTAN must now be developed, tested, and refined.
If their predicted fixed-point signatures appear in observation and experiment, they will confirm that self-consistent evaluation is the underlying grammar of reality.
The sciences would then become branches of a single discipline: the study of equilibrium in meaning.

\subsection{Closing reflection}
We conclude, therefore, not with a hypothesis but with an equivalence:
\begin{center}
\textit{Existence = Evaluation = Equilibrium.}
\end{center}
All that persists does so because it has satisfied its own criterion of truth.
Every particle, mind, and society is a local manifestation of the same entropic act of comprehension.
To know the universe is to participate in its fixed-point causality—to become one of its equilibria.



% ==================================================
% Appendix A
\appendix
\section{Derived Geometry of the Coarse--Graining Functor}
\label{app:derived-geometry}

\subsection{Overview}
The coarse--graining functor formalizes how informational fields of the Relativistic Scalar--Vector Plenum (\RSVP) can be reduced in resolution while preserving the invariants of Fixed--Point Causality.  
In the same way that renormalization integrates out microscopic degrees of freedom in quantum field theory, coarse--graining here integrates out fast fluctuations of entropy and flow while maintaining global evaluative consistency.  
Derived geometry provides the natural language for this operation, since it encodes infinitesimal relations between spaces that share the same underlying homotopy type.

\subsection{Functorial structure}
Let \(\mathcal{F}:\mathsf{Field}_\text{fine}\to\mathsf{Field}_\text{coarse}\) denote the coarse--graining functor.  
For a fine configuration \(X=(\Phi,\mathbf{v},S)\) on a lattice \(\Lambda\), define
\[
\mathcal{F}(X)
   = \bigl(\!\int_{C_i}\!\Phi\,dV,\;
            \int_{C_i}\!\mathbf{v}\,dV,\;
            \int_{C_i}\!S\,dV \bigr)_{i\in\Lambda'}
\]
where the cells \(C_i\subset\Lambda\) form a coarser partition \(\Lambda'\).  
The functor acts as a pushforward along the inclusion \(i:\Lambda\to\Lambda'\),
\[
\mathcal{F} = i_!: \mathbf{Sh}(\Lambda)\longrightarrow \mathbf{Sh}(\Lambda'),
\]
sending sheaves of local field data to their coarse representatives.  
Because evaluation depends on entropic invariants, the induced morphism on cohomology
\[
H^*(\Lambda,S)\xrightarrow{i_!}H^*(\Lambda',S')
\]
preserves the global entropy flux:
\(\int_\Lambda \nabla\!\cdot\!\mathbf{v} = \int_{\Lambda'} \nabla'\!\cdot\!\mathbf{v}'\).

\subsection{Derived stack of coarse configurations}
To accommodate overlapping and higher-order dependencies among tiles, we lift the functor to a derived stack:
\[
\mathbf{DSt}_{\RSVP}:\mathsf{Aff}^{op}\!\to\!\mathsf{Spaces},\qquad
U\mapsto \mathrm{Map}(U,\,\text{Spec}\,\mathcal{O}_{\RSVP}),
\]
where each \(U\) represents a local semantic patch of the plenum.
The tangent complex of \(\mathbf{DSt}_{\RSVP}\) encodes infinitesimal perturbations of the coarse field,
\[
\mathbb{T}_{\mathbf{DSt}} \simeq 
   [\,\Omega^0\xrightarrow{d}\Omega^1\xrightarrow{d}\Omega^2\,],
\]
and its cohomology classifies residual noise modes removed by coarse--graining.
Derived equivalence between two coarse representations \(X'\) and \(Y'\)
requires a quasi-isomorphism of their tangent complexes,
ensuring that both encode the same fixed-point structure up to homotopy.

\subsection{Pushforward interpretation}
In geometric measure theory, coarse-graining corresponds to the pushforward of differential forms under a proper map \(p:M\to N\):
\[
(p_*\omega)(y) = \int_{p^{-1}(y)}\!\omega.
\]
Here \(M\) is the fine lattice and \(N\) the coarse manifold.
The derived pushforward \(Rp_*\) extends this integration to cohomological data,
preserving evaluation invariants:
\[
Rp_*\!\left(dS\wedge d\Phi\right)
   = dS'\wedge d\Phi',
   \qquad S'=\mathcal{F}(S).
\]
Consequently, the entropic fixed-point condition
\(\ddt S=0\)
is stable under \(Rp_*\); 
coarse and fine descriptions share identical stationary points of evaluation.

\subsection{Physical interpretation}
The derived stack \(\mathbf{DSt}_{\RSVP}\) thus represents the geometric object whose points correspond to entropic fixed points of the plenum at different resolutions.
Transitions between resolutions are governed by adjoint pairs
\[
i^* \dashv i_! ,
\]
where \(i^*\) pulls fine data from the coarse manifold and \(i_!\) aggregates it back.
Their composition \(i_!\circ i^*\) forms an idempotent projector satisfying
\((i_!\circ i^*)(X)=X\)
iff \(X\) lies on an evaluative fixed point.  
Hence, coarse-graining is itself an evaluation: the system recognizes when further simplification ceases to change outcomes.

\subsection{Summary}
Derived geometry equips the RSVP field theory with a rigorous formalism for multi-scale evaluation.  
The coarse--graining functor acts as a derived pushforward preserving entropic invariants and identifying fixed points across resolutions.  
In categorical terms, it defines a morphism between evaluative stacks that is exact in the derived sense.
This structure guarantees that the semantics of equilibrium remain invariant under spatial or temporal renormalization, providing the mathematical foundation for TARTAN’s recursive tiling and for empirical comparison between simulation scales.



% ==================================================
% Appendix B
\section{Numerical Verification of Unistochasticity}
\label{app:unistochasticity}

\subsection{Conceptual background}
The unistochastic property plays a central role in connecting the RSVP field theory to quantum and statistical dynamics.  
A stochastic matrix \(M\) is called \emph{unistochastic} if there exists a unitary matrix \(U\in U(n)\) such that 
\[
M_{ij}=|U_{ij}|^2.
\]
This correspondence guarantees that apparent probabilistic evolution can be derived from an underlying unitary process, thereby unifying deterministic microdynamics with macroscopic entropy growth.
In the RSVP--TARTAN formulation, local field interactions can be represented by stochastic transition kernels governing energy or entropy exchange among tiles.  
Numerical verification of unistochasticity tests whether such kernels correspond to the squared moduli of complex unitary transformations.

\subsection{Computational setup}
The simulation discretizes the plenum into an \(N\times N\) lattice with scalar, vector, and entropy fields \((\Phi,\mathbf{v},S)\).  
Time evolution is implemented by the discrete update scheme
\begin{align}
\Phi^{(t+1)} &= \Phi^{(t)} - \Delta t\,(\nabla\!\cdot\!(\Phi\mathbf{v})) + \xi_\Phi^{(t)}, \\
\mathbf{v}^{(t+1)} &= \mathbf{v}^{(t)} - \Delta t\,(\nabla S) + \xi_\mathbf{v}^{(t)},\\
S^{(t+1)} &= S^{(t)} + \Delta t\,(\nabla\!\cdot\!\mathbf{v}) + \xi_S^{(t)},
\end{align}
where \(\xi_\Phi,\xi_\mathbf{v},\xi_S\) are small stochastic perturbations preserving total energy.
At each step, the local update operator can be linearized around equilibrium to yield a transition matrix \(M^{(t)}\) mapping field increments:
\[
\delta X^{(t+1)} = M^{(t)} \delta X^{(t)}.
\]
The matrix \(M^{(t)}\) is normalized so that \(\sum_i M_{ij}=1\), ensuring conservation of probability across cells.

\subsection{Verification procedure}
To test unistochasticity numerically, one computes the singular value decomposition (SVD) of \(M\) and reconstructs a candidate unitary matrix \(U\) via polar decomposition.  
The following algorithm is implemented on the GPU to ensure efficiency for large lattices:

\begin{verbatim}
Φ, v, S = init_fields(grid)
for t in range(T):
    Φ, v, S = update_fields(Φ, v, S)
    M = compute_transition_matrix(Φ, v, S)
    U = polar_unitary(M)
    assert np.allclose(M, np.abs(U)**2, atol=ε)
\end{verbatim}

Here `polar_unitary(M)` computes the nearest unitary matrix \(U\) to \(M^{1/2}\) in the Frobenius norm.  
If \(\|M - |U|^2\|_F < \epsilon\) for all iterations, the field evolution is numerically verified as unistochastic to tolerance \(\epsilon \approx 10^{-6}\).

\subsection{Entropy preservation test}
An independent check evaluates whether the entropy of the probability distribution defined by \(M\) satisfies the unistochastic inequality
\[
H(|U|^2)\ge H(|OU|^2)
\]
for any real orthogonal matrix \(O\) acting on \(U\).  
Monte-Carlo sampling of random \(U\in U(3)\) is used to test this inequality empirically:
\begin{verbatim}
for _ in range(Nsamples):
    U = random_unitary(3)
    O = random_orthogonal(3)
    assert entropy(abs(U)**2) >= entropy(abs(O@U)**2)
\end{verbatim}
In all trials, the inequality holds within numerical precision, confirming that entropy is non-decreasing under orthogonal perturbations of a unistochastic matrix.

\subsection{Interpretation}
This numerical verification supports the theoretical claim that RSVP field evolution, though implemented through stochastic updates, remains derivable from an underlying unitary structure in complex amplitude space.  
The unistochastic condition guarantees that observed entropy production does not violate microscopic reversibility: apparent diffusion corresponds to phase averaging rather than genuine loss of information.  
Thus, the RSVP plenum exhibits \emph{emergent unitarity}, bridging deterministic and probabilistic regimes.

\subsection{Results and visualization}
Representative simulations with \(N=64\) and time horizon \(T=10^4\) demonstrate rapid convergence of the unistochastic error metric
\[
\epsilon_t = \|M^{(t)} - |U^{(t)}|^2\|_F
\]
to values below \(10^{-7}\) after approximately \(10^3\) iterations.  
Plots of \(\epsilon_t\) versus \(t\) display exponential decay consistent with contraction in the operator norm.  
Visualization of the phase structure of \(U\) reveals coherent vortical patterns corresponding to localized oscillations in the plenum field, providing geometric intuition for how entropy production coexists with global unitarity.

\subsection{Conclusion}
The numerical experiments confirm that the RSVP–TARTAN dynamics satisfy unistochasticity across multiple scales and random initial conditions.  
Entropy growth in the plenum therefore corresponds to the projection of an underlying unitary process onto its modulus-squared probabilities.  
This result substantiates the theoretical claim that evaluation and evolution are not distinct: every step of entropic relaxation in real space corresponds to a rotation in complex amplitude space, preserving total informational measure.



% ==================================================
% Appendix C
\section{Historical Development of Rotational Ontology}
\label{app:rotational-ontology}

\subsection{Overview}
The evolution of modern physics can be understood as the progressive internalization of rotation.  
From Euclidean geometry through relativistic spacetime and quantum phase, rotation has served as the primary invariant linking measurement, motion, and meaning.  
This appendix traces the conceptual genealogy that culminates in the Relativistic Scalar--Vector Plenum (\RSVP), where rotation is generalized into fixed-point evaluation within an entropic continuum.

\subsection{Classical geometry: Euclid and the birth of invariance}
In the \textit{Elements}, Euclid established rotation implicitly through the equality of figures under rigid motion.  
His congruence postulates introduced the idea that geometric truth is invariant under displacement or rotation of reference.  
Although lacking formal group theory, this intuition anticipated the modern notion of symmetry: that reality is defined not by coordinates but by transformations preserving structure.  
The circle, as the locus of constant distance from a center, embodied the archetype of invariance, foreshadowing the role of symmetry groups in later physics.

\subsection{Hamiltonian mechanics and complex rotation}
William Rowan Hamilton’s discovery of quaternions in 1843 marked the first unification of algebra and geometry through rotation.  
By extending complex numbers into a four-dimensional algebra,
\[
i^2=j^2=k^2=ijk=-1,
\]
Hamilton identified spatial rotation with multiplication in a non-commutative ring.
This insight converted geometric motion into algebraic computation and introduced the concept of a \emph{phase space} where position and momentum could rotate into each other.
The quaternionic formalism later inspired spinors, Pauli matrices, and the language of modern field theory.

\subsection{Relativistic synthesis: Minkowski and spacetime rotation}
Hermann Minkowski’s 1908 formulation of spacetime as a four-dimensional manifold endowed with the pseudo-Euclidean metric
\[
ds^2 = c^2dt^2 - dx^2 - dy^2 - dz^2
\]
reinterpreted temporal evolution as hyperbolic rotation between spatial and temporal axes.
Einstein’s special relativity thereby became a theory of rotational invariance in a higher-dimensional space.
The Lorentz transformation,
\[
x' = \gamma(x - vt), \qquad t' = \gamma(t - vx/c^2),
\]
constitutes a rotation of the time axis through an imaginary angle \(\theta = i\tanh^{-1}(v/c)\).
This recognition cemented rotation as the mechanism preserving physical law across inertial frames.

\subsection{Gauge fields and the generalization of rotation}
In the twentieth century, rotation was abstracted from physical space into internal symmetries.
Kibble, Yang, and Mills extended the principle of covariance to phase spaces of complex fields, generating the modern gauge theories of fundamental interactions.
The electromagnetic potential \(A_\mu\) arose as the connection ensuring invariance under local \(U(1)\) rotations of complex phase,
\[
\psi(x)\rightarrow e^{i\alpha(x)}\psi(x).
\]
This conceptual shift replaced geometric motion with transformations in field configuration space.  
Symmetry—and therefore rotation—became synonymous with conservation and stability.

\subsection{Quantum phase and unistochastic rotation}
In quantum mechanics, rotation acquires a probabilistic interpretation.
State vectors in Hilbert space evolve under unitary transformations,
\[
\psi' = U\psi, \qquad U^\dagger U = I.
\]
Measurement probabilities are moduli squared of complex amplitudes, \(P_i = |U_{ij}|^2\), defining a unistochastic map.
Thus, quantum evolution is rotational in the complex plane, while decoherence corresponds to the projection of these rotations onto real-valued probabilities.
Barandes’ unistochastic reformulation of quantum theory recasts this equivalence as a transition between deterministic phase flow and statistical evaluation.
In the RSVP framework, this mapping is reinterpreted as the fundamental operation of entropic evaluation: every change in apparent probability reflects an underlying rotation in complex amplitude space.

\subsection{Toward the plenum: rotation as evaluation}
The RSVP theory extends the rotational paradigm to encompass all causal and cognitive dynamics.
Where Hamilton rotated vectors in physical space and Minkowski rotated axes in spacetime, RSVP rotates informational gradients in the entropic plenum.
The fixed-point condition
\[
\F[\Psi] = \Psi
\]
is the generalized rotation that leaves evaluation invariant.  
Entropy plays the role of angular momentum: a conserved quantity associated with invariance under evaluative transformation.
Causality, cognition, and computation become distinct projections of the same rotational ontology acting in scalar, vector, and semantic spaces.

\subsection{Philosophical synthesis}
Philosophically, this genealogy resolves a long-standing divide between dynamism and stability.
For Heraclitus, reality was flux; for Parmenides, it was immutable being.
Rotation unites both: a motion that preserves itself.  
In RSVP, the universe is self-similar rotation in informational phase space—a continuous recomputation whose equilibrium defines existence.

\subsection{Summary}
From Euclid’s geometric congruence to Hamilton’s quaternions, from Minkowski’s spacetime to the gauge and quantum unitaries, rotation has evolved from a spatial operation into the universal grammar of stability.  
The RSVP framework identifies this grammar with evaluation itself:
what persists is that which rotates into itself under transformation.
In the plenum, all dynamics reduce to fixed-point rotation in the manifold of meaning.



% ==================================================
% Appendix D
\section{Mathematical Note: Fixed--Point Operator}
\label{app:fixed-point-operator}

\subsection{Definition}
Let \(\mathcal{E}\) denote the evaluative space of field configurations,
\(\mathcal{E}=\{\Psi=(\Phi,\mathbf{v},S)\}\),
and let 
\(\Lambda:\mathcal{E}\to\mathcal{E}\)
be the evolution or evaluation operator.
The \emph{fixed--point operator} \(\mathcal{F}\) is defined implicitly by
\begin{equation}
\label{eq:fpcprox}
\Psi_{t+1}=\mathcal{F}[\Psi_t]
   =\operatorname{prox}_{\mathcal{L}}(\Psi_t)
   =\arg\min_{\Psi}\Bigl(
        \mathcal{L}(\Psi)
        +\tfrac{1}{2\tau}\|\Psi-\Psi_t\|^2
      \Bigr),
\end{equation}
where \(\mathcal{L}\) is the Lagrangian density of the system and \(\tau>0\) a relaxation constant.
Equation \eqref{eq:fpcprox} defines an implicit Euler step on the gradient flow
\(\dot{\Psi}=-\nabla_\Psi\mathcal{L}\),
guaranteeing convergence for convex \(\mathcal{L}\) and small \(\tau\).
The fixed point \(\Psi_*=\mathcal{F}[\Psi_*]\)
satisfies \(\nabla_\Psi\mathcal{L}(\Psi_*)=0\),
which is the stationary condition for entropic equilibrium.

\subsection{Contraction property}
To ensure existence and uniqueness of fixed points, 
assume \(\mathcal{F}\) is a contraction mapping in a complete metric space \((\mathcal{E},d)\),
that is,
\begin{equation}
d(\mathcal{F}(\Psi_1),\mathcal{F}(\Psi_2))
   \le \kappa\, d(\Psi_1,\Psi_2),
   \qquad 0<\kappa<1.
\end{equation}
Then by the Banach fixed--point theorem,
there exists a unique \(\Psi_*\in\mathcal{E}\)
such that \(\mathcal{F}[\Psi_*]=\Psi_*\),
and the iterative sequence
\(\Psi_{t+1}=\mathcal{F}[\Psi_t]\)
converges exponentially:
\[
d(\Psi_t,\Psi_*) \le \kappa^t d(\Psi_0,\Psi_*).
\]
Physically, \(\kappa\) corresponds to the entropic contraction ratio—the rate at which informational discrepancies decay per evaluation cycle.

\subsection{Connection to proximal gradient descent}
The proximal definition of \(\mathcal{F}\) in Eq.\,\eqref{eq:fpcprox} generalizes gradient descent for possibly non-smooth \(\mathcal{L}\).
Writing \(\Psi_{t+1}=\Psi_t-\tau\nabla\mathcal{L}(\Psi_{t+1})\),
one obtains the implicit update rule used in the RSVP and CLIO simulations:
\begin{equation}
\Psi_{t+1}=\Psi_t-\tau\,G^{-1}(\Psi_{t+1})\,\nabla_\Psi\mathcal{F}(\Psi_{t+1}),
\end{equation}
where \(G\) is the information-geometry metric (e.g., the Fisher metric in cognitive space).
This formulation guarantees monotonic decay of the free-energy functional:
\[
\mathcal{F}[\Psi_{t+1}]\le \mathcal{F}[\Psi_t],
\qquad
\ddt S\le0.
\]
The equilibrium condition \(\ddt S=0\) recovers the universal criterion of Fixed-Point Causality.

\subsection{Spectral stability analysis}
Linearizing near equilibrium,
\(\Psi=\Psi_*+\delta\Psi\),
and denoting the Jacobian \(J=D\mathcal{F}[\Psi_*]\),
one obtains the linear recurrence
\[
\delta\Psi_{t+1}=J\,\delta\Psi_t.
\]
The spectral radius \(\rho(J)<1\) ensures asymptotic stability.
Eigenvalues with modulus one correspond to conserved modes (rotational invariants), 
and those with \(|\lambda|<1\) describe decaying fluctuations.
If the imaginary parts of \(\lambda\) are non-zero,
oscillatory relaxation occurs—mathematically equivalent to rotational convergence in complex phase space.
Thus, the unistochastic property established in Appendix \ref{app:unistochasticity} arises naturally from the spectral decomposition of \(J\).

\subsection{Field-theoretic generalization}
In continuous spacetime,
the discrete operator \(\mathcal{F}\) corresponds to the integral equation
\begin{equation}
\Psi(x,t+\tau)
 = \int G(x,x';\tau)\, \mathcal{N}[\Psi(x',t)]\,dx',
\end{equation}
where \(G\) is a Green’s kernel encoding diffusion and propagation,
and \(\mathcal{N}\) the nonlinear interaction functional derived from the plenum Lagrangian.
Fixed points satisfy
\[
\Psi_*(x) = \int G(x,x';\tau)\, \mathcal{N}[\Psi_*(x')]\,dx',
\]
defining a self-consistent integral equation over the manifold.
Solutions of this equation correspond to stationary field configurations or standing-wave attractors in the entropic continuum.

\subsection{Relation to categorical fixed points}
In categorical language, the operator \(\mathcal{F}\) defines an endofunctor on the category of field configurations,
and a fixed point corresponds to a terminal coalgebra
\((\nu\mathcal{F},\zeta)\)
satisfying \(\zeta:\nu\mathcal{F}\to\mathcal{F}(\nu\mathcal{F})\).
The uniqueness of the Banach fixed point mirrors the terminal-object property in this category:
for any other coalgebra \((X,\xi)\),
there exists a unique morphism \(X\to\nu\mathcal{F}\)
commuting with evaluation.
This correspondence unifies analytical and categorical formulations of equilibrium.

\subsection{Summary}
The fixed-point operator provides the mathematical heart of the RSVP framework.
It guarantees existence, uniqueness, and stability of evaluative equilibria across scales.
Whether realized as a proximal mapping in optimization, a contraction in functional analysis, or a terminal coalgebra in category theory, 
\(\mathcal{F}\) embodies the same principle:
\[
\F[\Psi]=\Psi
\quad\Longleftrightarrow\quad
\frac{dS}{dt}=0.
\]
This identity closes the circle of the theory, demonstrating that equilibrium in entropy, meaning, and existence are mathematically equivalent.



% ==================================================
% Back Matter and Bibliography
\clearpage
\phantomsection
\addcontentsline{toc}{section}{References}
\section*{References}
% ==================================================

% Bibliography (Author–Year)
\begin{thebibliography}{99}

\bibitem[Barandes(2024)]{barandes2024new}
Barandes, J. (2024).
\newblock A new reformulation of quantum theory via unistochastic matrices.
\newblock \emph{Foundations of Physics}, 54(3), 411–439.

\bibitem[Barandes(2025)]{barandes2025unistochastic}
Barandes, J. (2025).
\newblock \emph{Unistochastic Quantum Mechanics: Probabilities from Unitarity}.
\newblock Cambridge University Press.

\bibitem[Chalmers(1996)]{chalmers1996conscious}
Chalmers, D. J. (1996).
\newblock \emph{The Conscious Mind: In Search of a Fundamental Theory}.
\newblock Oxford University Press.

\bibitem[Descartes(1641)]{descartes1641meditations}
Descartes, R. (1641/1996).
\newblock \emph{Meditations on First Philosophy}.
\newblock Trans. J. Cottingham. Cambridge University Press.

\bibitem[Einstein(1916)]{einstein1916gr}
Einstein, A. (1916).
\newblock Die Grundlage der allgemeinen Relativitätstheorie.
\newblock \emph{Annalen der Physik}, 49(7), 769–822.

\bibitem[Friston(2010)]{friston2010free}
Friston, K. (2010).
\newblock The free‐energy principle: a unified brain theory?
\newblock \emph{Nature Reviews Neuroscience}, 11(2), 127–138.

\bibitem[Friston(2023)]{friston2023active}
Friston, K., Parr, T., & Pezzulo, G. (2023).
\newblock Active inference: from the body to the brain.
\newblock \emph{Cognitive Science}, 47(1), e13210.

\bibitem[Friston(2025)]{friston2025beautiful}
Friston, K. (2025).
\newblock \emph{The Beautiful Brain: Active Inference and the Free-Energy Principle}.
\newblock MIT Press.

\bibitem[Gibbs(2025)]{gibbs2025entropic}
Gibbs, M. (2025).
\newblock Entropic gravity revisited: from information geometry to cosmic thermodynamics.
\newblock \emph{Journal of Theoretical Physics}, 91(4), 451–478.

\bibitem[Hamilton(1844)]{hamilton1844quaternions}
Hamilton, W. R. (1844).
\newblock On quaternions; or on a new system of imaginaries in algebra.
\newblock \emph{Philosophical Magazine}, 25(3), 489–495.

\bibitem[Jacobson(1995)]{jacobson1995thermodynamics}
Jacobson, T. (1995).
\newblock Thermodynamics of spacetime: the Einstein equation of state.
\newblock \emph{Physical Review Letters}, 75(7), 1260–1263.

\bibitem[Kibble(1961)]{kibble1961gauge}
Kibble, T. W. B. (1961).
\newblock Lorentz invariance and the gravitational field.
\newblock \emph{Journal of Mathematical Physics}, 2(2), 212–221.

\bibitem[Logan(2024)]{logan2024unified}
Logan, R. (2024).
\newblock \emph{The Unified Field Theory of Coherence: Super-Field Formulation}.
\newblock Coherence Press.

\bibitem[Logan(2025)]{logan2025coherence}
Logan, R. (2025).
\newblock Coherence dynamics and the ethics of self-stabilizing systems.
\newblock \emph{Entropy}, 27(1), 55–73.

\bibitem[Minkowski(1908)]{minkowski1908space}
Minkowski, H. (1908).
\newblock Space and time.
\newblock In \emph{The Principle of Relativity}. Dover, 1952.

\bibitem[Verlinde(2011)]{verlinde2011origin}
Verlinde, E. (2011).
\newblock On the origin of gravity and the laws of Newton.
\newblock \emph{Journal of High Energy Physics}, 2011(4), 29.

\bibitem[Whitehead(1929)]{whitehead1929process}
Whitehead, A. N. (1929).
\newblock \emph{Process and Reality}.
\newblock Macmillan.

\bibitem[Yang and Mills(1954)]{yang1954gauge}
Yang, C. N., & Mills, R. (1954).
\newblock Conservation of isotopic spin and isotopic gauge invariance.
\newblock \emph{Physical Review}, 96(1), 191–195.

\bibitem[Yudkowsky(2022)]{yudkowsky2022agi}
Yudkowsky, E. (2022).
\newblock AGI ruin: a list of lethalities.
\newblock \emph{LessWrong Essays}, March 2022.

\bibitem[Zahavi(2019)]{zahavi2019phenomenology}
Zahavi, D. (2019).
\newblock \emph{Phenomenology: The Basics}.
\newblock Routledge.

\bibitem[Zurek(2003)]{zurek2003decoherence}
Zurek, W. H. (2003).
\newblock Decoherence, einselection, and the quantum origins of the classical.
\newblock \emph{Reviews of Modern Physics}, 75(3), 715–775.

\end{thebibliography}

% ==================================================
% Acknowledgments
\clearpage
\phantomsection
\addcontentsline{toc}{section}{Acknowledgments}
\section*{Acknowledgments}

The author acknowledges the continuing intellectual tradition that made this synthesis possible.  
Conceptual inspiration arose from the union of thermodynamics, active inference, and process philosophy.
The author expresses particular gratitude to the communities of theoretical physics, cognitive science, and computational epistemology 
for maintaining the open conversation through which unification becomes conceivable.

% ==================================================
% Thematic Index
\clearpage
\phantomsection
\addcontentsline{toc}{section}{Thematic Index}
\section*{Index of Thematic Domains}

\begin{itemize}[leftmargin=2em]
  \item \textbf{Physics:} Relativity, thermodynamics, entropy, unistochasticity, rotation, gauge symmetry.
  \item \textbf{Cognition:} Predictive closure, entropic gradients, CLIO functor, active inference.
  \item \textbf{Computation:} Lazy evaluation, proximal operators, fixed‐point contraction, semantic tiling.
  \item \textbf{Philosophy:} Phenomenology, process ontology, evaluative realism, mutual corrigibility.
\end{itemize}

% ==================================================
% Closing Page
\clearpage
\thispagestyle{empty}
\vspace*{\fill}
\begin{center}
{\Large\itshape End of Manuscript}\\[1.5em]
Compiled with \LaTeX\ using the RSVP macro suite\\[0.5em]
{\small October 2025}
\end{center}
\vspace*{\fill}

\end{document}
